% Appendices. main.tex inputs this file after the references, following
% \startappendices. Sections are lettered (Appendix A, B, ...); a \ref to any
% label here returns the letter alone, so the text cites them as
% Appendix~\ref{app:theory}. The files input below start at subsection level.
%
%   A  Proofs of the analytic results   appendix_theory (with app:gradnorm)
%   B  Experimental details             appendix_setup, appendix_scoring
%   C  Pricing and comparison tables    generated tables moved from the text
%   D  Supplementary results            supp_audit, supp_mechanism,
%                                       supp_confound, supp_generality
% Tables and figures are numbered per appendix (Table C.1, Fig. D.1), as
% \startappendices in main.tex arranges.

% Most generated tables here are about two thirds of a page tall: too tall to
% share a page with text, and below the main text's float-page threshold
% (0.80, main.tex), so they would wait for the next \FloatBarrier. At 0.70
% each forms a float page of its own as soon as it is reached.
\renewcommand{\floatpagefraction}{0.70}

\section{Proofs of the analytic results}\label{app:theory}
\IfFileExists{sections/appendix_theory.tex}{% Propositions and proofs cited from Sections~\ref{sec:background}, \ref{sec:family}
% and \ref{sec:confound}. Input by appendix.tex under the section heading that
% carries \label{app:theory}, so this file starts at subsection level. Uses the
% `proposition' environment and amsthm's `proof', both set up in main.tex.

\subsection{Setting}
\label{app:theory-setting}

Fix a timestep $t$ and write $\sigma = \sigma_t$. We treat one channel; with
several channels every quantity below is averaged over channels, which leaves
the statements unchanged. The residual is $r \in \mathbb{R}^N$. A linear map
$W : \mathbb{R}^N \to \mathbb{R}^{N'}$ has its coefficients partitioned into $J$
bands $b \in \mathcal{B}$, with $P_b$ the restriction to the $n_b$ coefficients
of band $b$ and $\sum_b n_b = N'$. A complex unitary transform such as the DFT
is covered by the same arguments with the adjoint $W^{*}$ in place of
$W^\top$. For non-negative weights $\omega_b(\sigma)$
the per-sample loss of Eq.~\eqref{eq:family} and the pixel mean square are
\begin{equation}
  \ell_\sigma(r) = \sum_{b} \frac{\omega_b(\sigma)}{n_b} \lVert P_b W r \rVert_2^2
                 = r^\top Q_\sigma r,
  \qquad
  Q_\sigma = W^\top D_\sigma W,
  \qquad
  m(r) = \frac{\lVert r \rVert_2^2}{N},
  \label{eq:app-loss}
\end{equation}
where $D_\sigma$ is diagonal with entry $\omega_b(\sigma)/n_b$ on the
coefficients of band $b$. For $r \neq 0$ let
$\rho_b(r) = \operatorname{ms}(P_b W r)/m(r)$ be the mean square of band $b$
relative to the pixel mean square. If $W^\top W = I$, Eq.~\eqref{eq:parseval}
gives $\sum_b (n_b/N)\,\rho_b(r) = 1$; for orthogonal $W$ ($N' = N$),
$\rho_b \equiv 1$ when the residual's
energy is spread over the coefficients in proportion to band size (as for a
white residual in expectation).

\subsection{Band weighting leaves the minimiser unchanged}
\label{app:minimiser}

\begin{proposition}[Population minimiser]
\label{prop:minimiser}
Let $W$ be injective, and let $\omega_b(\sigma_t) > 0$ for every band and
timestep, with the weights depending on $(x_0, \epsilon, t)$ only through $t$.
For a measurable predictor $f$ let
$\mathcal{R}(f) = \mathbb{E}_{x_0,\epsilon,t}\big[\ell_{\sigma_t}\big(f(x_t,t) - \epsilon\big)\big]$.
Then $\mathcal{R}(f) \ge \mathcal{R}(f^\star)$ with
$f^\star(x_t, t) = \mathbb{E}[\epsilon \mid x_t, t]$, and equality holds only
if $f = f^\star$ almost everywhere. In particular, every such weighting,
including the uniform one that gives $\mathcal{L}_{\mathrm{simple}}$, has the
same minimiser.
\end{proposition}

\begin{proof}
For $r \neq 0$,
$r^\top Q_\sigma r = \sum_b (\omega_b/n_b)\lVert P_b W r\rVert_2^2
\ge (\min_b \omega_b/n_b)\,\lVert W r \rVert_2^2 > 0$, because $W r \neq 0$;
so $Q_\sigma$ is positive definite. Fix $(x_t, t)$ and write
$\mu = f^\star(x_t, t)$ and $y = f(x_t, t)$. Since $Q_{\sigma_t}$, $y$ and
$\mu$ are functions of $(x_t, t)$ and
$\mathbb{E}[\epsilon - \mu \mid x_t, t] = 0$, the cross term vanishes and
\begin{equation*}
  \mathbb{E}\big[(y-\epsilon)^\top Q_{\sigma_t} (y-\epsilon) \,\big|\, x_t, t\big]
  = (y-\mu)^\top Q_{\sigma_t} (y-\mu)
  + \mathbb{E}\big[(\epsilon-\mu)^\top Q_{\sigma_t} (\epsilon-\mu) \,\big|\, x_t, t\big].
\end{equation*}
The second term does not depend on $f$ and is finite because
$\mathbb{E}\lVert\epsilon\rVert^2 < \infty$; the first is non-negative and
vanishes only at $y = \mu$. Taking the expectation over $(x_t, t)$ completes
the proof.
\end{proof}

Three remarks delimit the result. (i) The argument is unchanged if $\ell$ is
multiplied by a positive timestep weight, so it covers the normalised
weightings of Section~\ref{sec:confound} and any product of a timestep
weighting with a band weighting. (ii) What the weighting does change is the
attained minimum (the second term above) and, when the model class does not
contain $f^\star$, which approximation is optimal: by the decomposition above,
the best predictor in a restricted class is the one closest to $f^\star$ in
the norm $\big(\mathbb{E}[(f - f^\star)^\top Q_{\sigma_t} (f - f^\star)]\big)^{1/2}$,
which depends on $\omega$. (iii) The proposition does not cover
weights that depend on $\epsilon$. If the weights are computed from the target,
as in the spatially adaptive arm of Section~\ref{sec:extensions}, whose map is
built from the target's coefficients, the loss is a quadratic form
$Q(\epsilon)$ and its minimiser is the weighted conditional mean
$\mathbb{E}[Q(\epsilon) \mid x_t, t]^{-1}\,\mathbb{E}[Q(\epsilon)\,\epsilon \mid x_t, t]$.
If the weights are computed from the residual and detached, as in FFL
(Eq.~\eqref{eq:ffl}), training does not descend a fixed objective; its
stationary points satisfy
$\mathbb{E}\big[F^{*}\operatorname{diag}(w(R))\,R \,\big|\, x_t, t\big] = 0$
with $R = F(f - \epsilon)$. Because $w$ is even in $R$, the conditional mean
satisfies this condition when the conditional law of $\epsilon$ is symmetric
about its mean, but not in general.

\subsection{The timestep component of a band weighting}
\label{app:confound}

\begin{proposition}[Timestep and shape decomposition]
\label{prop:confound}
Let $W$ be orthogonal ($N' = N$ and $W^\top W = I$) and let the weights
$\omega_b(\sigma) \ge 0$ depend on $\sigma$ only. Write $S(\sigma) = \sum_b \omega_b(\sigma)$ and
$\delta_b(\sigma) = \omega_b(\sigma) - (n_b/N)\,S(\sigma)$, so that
$\sum_b \delta_b = 0$. Then:
\begin{enumerate}
  \item[(a)] For every residual $r$,
  \begin{equation}
    \ell_\sigma(r)
    = S(\sigma)\, m(r)
    + \sum_b \delta_b(\sigma)\,\big[\operatorname{ms}(P_b W r) - m(r)\big]
    = m(r) \sum_b \omega_b(\sigma)\,\rho_b(r) .
    \label{eq:shape-decomp}
  \end{equation}
  The second term vanishes identically when $\omega_b \propto n_b$, and for a
  given residual whenever $\rho_b(r) = 1$ for every band with
  $\delta_b \neq 0$.
  \item[(b)] Let
  $\bar\rho_b(t) = \mathbb{E}[\operatorname{ms}(P_b W r) \mid t]/\mathbb{E}[m(r) \mid t]$.
  Then $\mathbb{E}[\ell_{\sigma_t}(r) \mid t] = \Lambda(t)\,\mathbb{E}[m(r) \mid t]$ with
  \begin{equation}
    \Lambda(t) = \sum_b \omega_b(\sigma_t)\,\bar\rho_b(t)
    = S(\sigma_t) + \sum_b \delta_b(\sigma_t)\,\big(\bar\rho_b(t) - 1\big)
    = \sum_b \frac{N}{n_b}\,\omega_b(\sigma_t)\, e_b(t),
    \label{eq:eff-weight}
  \end{equation}
  where $e_b(t) = (n_b/N)\,\bar\rho_b(t)$ are the residual's band energy
  fractions at timestep $t$ ($\sum_b e_b = 1$). The effective timestep weight
  $\Lambda(t)$ is thus a convex combination of the pixel-equivalent weights
  $(N/n_b)\,\omega_b(\sigma_t)$, and lies between their minimum and maximum.
  \item[(c)] For AWWL (Eq.~\eqref{eq:audited}; $J = 4$, $n_b = N/4$),
  $S(\sigma) = (1-\alpha) + (2\alpha-1)\lambda_p(\sigma)$ and
  \begin{align}
    \ell_\sigma(r) &= S(\sigma)\, m(r)
      + \Big(w_{\mathrm{LL}}(\sigma) - \frac{w_{\det}(\sigma)}{3}\Big)
        \big[\operatorname{ms}(r_{\mathrm{LL}}) - m(r)\big],
    \label{eq:awwl-decomp}\\
    \Lambda(t) &= \frac{4}{3}\,w_{\det}(\sigma_t)
      + \Big(w_{\mathrm{LL}}(\sigma_t) - \frac{w_{\det}(\sigma_t)}{3}\Big)\,\bar\rho_{\mathrm{LL}}(t),
    \label{eq:awwl-eff}
  \end{align}
  with $\bar\rho_{\mathrm{LL}}(t) \in [0, 4]$. The residual-independent weight
  $S$ is constant in $\sigma$ if and only if $\alpha = 1/2$ (for $p > 0$). For
  $\alpha \neq 1/2$ and any function $\phi$ of the timestep, the correlation over
  timesteps satisfies
  $\operatorname{corr}_t\big(S(\sigma_t), \phi(t)\big)
  = \operatorname{sign}(2\alpha - 1)\,\operatorname{corr}_t\big(\lambda_p(\sigma_t), \phi(t)\big)$,
  independently of $\lvert\alpha - 1/2\rvert$.
  \item[(d)] Rescaling the AWWL weights to a constant total $c$, that is,
  replacing $\omega_b$ by $\tilde\omega_b = c\,\omega_b/S$, removes the
  residual-independent timestep weight exactly and leaves
  $\tilde\Lambda(t) = c + c\,\big(w_{\mathrm{LL}} - w_{\det}/3\big)\big(\bar\rho_{\mathrm{LL}}(t) - 1\big)/S$,
  evaluated at $\sigma_t$. It equals $c$ at every timestep if and only if, at
  every timestep, either the approximation band carries a quarter of the
  residual's energy ($\bar\rho_{\mathrm{LL}}(t) = 1$) or the four per-band
  weights are equal ($w_{\mathrm{LL}} = w_{\det}/3$, which holds only at
  $\lambda_p(\sigma) = (1-\alpha)/(1+2\alpha)$).
\end{enumerate}
\end{proposition}

\begin{proof}
(a) Since $W^\top W = I$, Eq.~\eqref{eq:parseval} gives
$\sum_b n_b \operatorname{ms}(P_b W r) = \lVert r \rVert_2^2 = N m(r)$. Hence
\[
  \textstyle\sum_b \delta_b \operatorname{ms}(P_b W r)
  = \ell_\sigma(r) - (S/N)\sum_b n_b \operatorname{ms}(P_b W r)
  = \ell_\sigma(r) - S\,m(r),
\]
and $\sum_b \delta_b\, m(r) = 0$, which gives the first equality; the second is
the definition of $\rho_b$. If
$\omega_b \propto n_b$ then $\delta_b = 0$ for all $b$.
(b) Take conditional expectations in Eq.~\eqref{eq:shape-decomp} at timestep
$t$ and divide by $\mathbb{E}[m \mid t]$. The last form uses
$\omega_b \bar\rho_b = (N/n_b)\,\omega_b\, e_b$; $e_b \ge 0$ and
$\sum_b e_b = 1$ by Eq.~\eqref{eq:parseval}.
(c) With four bands of size $N/4$, $\sum_b \operatorname{ms}(r_b) = 4\,m(r)$, so
$\sum_{b \in \{\mathrm{LH},\mathrm{HL},\mathrm{HH}\}} \big(\operatorname{ms}(r_b) - m\big)
= -\big(\operatorname{ms}(r_{\mathrm{LL}}) - m\big)$. The three detail bands
share $\delta_{\det} = w_{\det}/3 - S/4$, and
$\delta_{\mathrm{LL}} - \delta_{\det} = w_{\mathrm{LL}} - w_{\det}/3$, so the
second term of Eq.~\eqref{eq:shape-decomp} equals
$(\delta_{\mathrm{LL}} - \delta_{\det})\big(\operatorname{ms}(r_{\mathrm{LL}}) - m\big)$.
Substituting $S = w_{\mathrm{LL}} + w_{\det}$ gives the expression for
$\Lambda$. The bound $\bar\rho_{\mathrm{LL}} \le 4$ follows from
$\lVert r_{\mathrm{LL}} \rVert^2 \le \lVert r \rVert^2$. $S$ is affine in
$\lambda_p$ with slope $2\alpha - 1$ by Eq.~\eqref{eq:total}, and $\lambda_p$
is strictly increasing on $(0,1)$ for $p > 0$; a Pearson correlation is
invariant under increasing affine maps of either argument and changes sign
under decreasing ones.
(d) Replace $\omega_b$ by $c\,\omega_b/S$ in (c).
\end{proof}

In Eq.~\eqref{eq:shape-decomp} the first term is the pixel MSE weighted by
$S(\sigma_t)$, a timestep weight that does not depend on the residual. The
second applies the zero-sum weights $\delta_b$ to the departures of the band
mean squares from the pixel mean square; it carries the frequency shape of the
weighting. By Eq.~\eqref{eq:family-x0}, the weight that the objective places on
the clean-image error at timestep $t$ is $\mathrm{SNR}_t\,\Lambda(t)$. Parts
(b) and (d) locate the part of the timestep weighting that no rescaling of the
weights by a function of $\sigma$ fixed in advance removes: the interaction of the zero-sum part $\delta_b$ with the
residual's band energy fractions, which change with $\sigma_t$
(Proposition~\ref{prop:residual}).

\subsection{What the optimiser receives}
\label{app:scale}

\begin{proposition}[Loss scale and gradient scale]
\label{prop:scale}
\begin{enumerate}
  \item[(a)] Consider the AdamW update
  \begin{equation}
    \theta_{k+1} = (1 - \eta_k\zeta)\,\theta_k
      - \eta_k\,\frac{\hat m_k}{\sqrt{\hat v_k} + \varepsilon_{\mathrm{A}}},
    \label{eq:adamw}
  \end{equation}
  where $\eta_k$ is the learning rate, $\zeta$ the decoupled weight decay,
  $\hat m_k$ and $\hat v_k$ the bias-corrected exponential moving averages of
  the gradients $g_k$ and of $g_k \odot g_k$, and the square root and division
  act elementwise. If every gradient is multiplied by a constant $c > 0$, the
  iterates are those of the unscaled run with $\varepsilon_{\mathrm{A}}$
  replaced by $\varepsilon_{\mathrm{A}}/c$. With gradient-norm clipping at a
  threshold $\tau$ applied before the update, the same holds if clipping is
  inactive on every step of both runs, and the two runs coincide exactly if it
  is active on every step of both.
  \item[(b)] For one sample at noise level $\sigma$,
  $\nabla_\theta \ell_\sigma = \sum_b \omega_b(\sigma)\, q_b$ with
  $q_b = \nabla_\theta \operatorname{ms}(P_b W r)$. If $\lVert q_b \rVert = G$ for
  every band, then $\lVert \nabla_\theta \ell_\sigma \rVert \le G\,S(\sigma)$, with
  equality when the $q_b$ are positively parallel. If in addition
  $\langle q_b, q_c \rangle \ge 0$ for all $b, c$, then
  $\lVert \nabla_\theta \ell_\sigma \rVert \ge G\,\big(\sum_b \omega_b(\sigma)^2\big)^{1/2}$,
  with equality when the $q_b$ are mutually orthogonal.
  \item[(c)] Let $W$ be orthogonal ($N' = N$, $W^\top W = W W^\top = I$). The
  output-space band gradients
  $h_b = \nabla_{\hat\epsilon} \operatorname{ms}(P_b W r) = (2/n_b)\, W^\top P_b^\top P_b W r$
  are mutually orthogonal, so
  \begin{equation}
    \lVert \nabla_{\hat\epsilon}\, \ell_\sigma \rVert_2^2
    = \sum_b \Big(\frac{2\,\omega_b(\sigma)}{n_b}\Big)^2 \lVert P_b W r \rVert_2^2 ,
    \label{eq:gradnorm}
  \end{equation}
  and for $J$ bands of equal size with $\rho_b(r) = 1$,
  $\lVert \nabla_{\hat\epsilon}\, \ell_\sigma \rVert_2 = (2\sqrt{J}/N)\,\lVert r \rVert_2\,\big(\sum_b \omega_b^2\big)^{1/2}$.
  For every $W$, the component of $\nabla_{\hat\epsilon}\, \ell_\sigma$ along
  the pixel-MSE gradient $\nabla_{\hat\epsilon}\, m$ has coefficient
  \begin{equation}
    \frac{\langle \nabla_{\hat\epsilon}\, \ell_\sigma, \nabla_{\hat\epsilon}\, m \rangle}
         {\lVert \nabla_{\hat\epsilon}\, m \rVert_2^2}
    = \frac{\ell_\sigma(r)}{m(r)} = \sum_b \omega_b(\sigma)\,\rho_b(r),
    \label{eq:mse-component}
  \end{equation}
  which equals $S(\sigma)$ when $\rho_b(r) = 1$ for every band.
  \item[(d)] For $\omega \in [0,\infty)^J$ with $S = \sum_b \omega_b$,
  \begin{equation}
    S^2/J \;\le\; \textstyle\sum_b \omega_b^2 \;\le\; S^2 ,
    \label{eq:cs-bounds}
  \end{equation}
  with equality on the left if and only if the weights are equal, and on the
  right if and only if at most one weight is non-zero. Holding $S$ fixed
  therefore leaves $\sum_b \omega_b^2$ free within a factor $J$, and holding
  $\sum_b \omega_b^2$ fixed leaves $S$ free within a factor $\sqrt{J}$.
\end{enumerate}
\end{proposition}

\begin{proof}
(a) $\hat m_k$ is linear in $g_1, \dots, g_k$ and $\hat v_k$ in
$g_1 \odot g_1, \dots, g_k \odot g_k$, so they scale by $c$ and $c^2$, and
$c\,\hat m_k/(c\sqrt{\hat v_k} + \varepsilon_{\mathrm{A}})
= \hat m_k/(\sqrt{\hat v_k} + \varepsilon_{\mathrm{A}}/c)$. The weight-decay term
does not involve the gradient, and the claim follows by induction on $k$.
Clipping multiplies $g$ by $\min(1, \tau/\lVert g \rVert)$; if it is inactive on
every step it changes nothing, and if it is active on every step both runs
receive $\tau\, g/\lVert g \rVert$, which does not depend on $c$.
(b) The upper bound is the triangle inequality. For the lower bound,
$\lVert \sum_b \omega_b q_b \rVert^2
= G^2 \sum_b \omega_b^2 + \sum_{b \neq c} \omega_b \omega_c \langle q_b, q_c \rangle
\ge G^2 \sum_b \omega_b^2$, since $\omega \ge 0$.
(c) $\nabla_{\hat\epsilon}\, \ell_\sigma = 2 Q_\sigma r = \sum_b \omega_b h_b$. For
$b \neq c$,
$\langle h_b, h_c \rangle \propto r^\top W^\top P_b^\top P_b\, W W^\top P_c^\top P_c W r = 0$,
because $W W^\top = I$ and $P_b P_c^\top = 0$; and
$\lVert h_b \rVert^2 = (4/n_b^2)\lVert P_b W r \rVert^2$, because $W^\top$ and
$P_b^\top$ preserve norms. With $n_b = N/J$ and
$\lVert P_b W r \rVert^2 = \lVert r \rVert^2/J$, Eq.~\eqref{eq:gradnorm} gives the
stated norm. Finally $\nabla_{\hat\epsilon}\, m = 2r/N$, and
$\langle 2 Q_\sigma r, 2r/N \rangle / \lVert 2r/N \rVert^2
= N\, r^\top Q_\sigma r / \lVert r \rVert^2 = \ell_\sigma(r)/m(r)$.
(d) The left inequality is the Cauchy--Schwarz inequality
$(\sum_b 1 \cdot \omega_b)^2 \le J \sum_b \omega_b^2$. The right one follows
from $S^2 = \sum_b \omega_b^2 + \sum_{b \neq c} \omega_b \omega_c$ with every cross
product non-negative, and it is an equality exactly when every cross product
vanishes.
\end{proof}

Part (c) is the orthogonal case of part (b): in output space the band
gradients are orthogonal, and with equal band energies the gradient norm scales
with $(\sum_b \omega_b^2)^{1/2}$. The parameter gradient is their image under
the transposed network Jacobian, which can align them; in the extreme case the
images are parallel and the norm scales with $S$. When the images have equal
norms and non-negative inner products (the assumptions of part (b)), the
per-timestep gradient contribution therefore lies between the two cases, and
the two normalisations,
$S$ held fixed and $\sum_b \omega_b^2$ held fixed, are the natural bounding
controls: the first fixes the contribution in the parallel case (and, in output
space, the component along the pixel-MSE gradient), the second fixes it in the
orthogonal case. By part (d) no rescaling fixes both unless the weights are
uniform, and both statements assume $\rho_b = 1$, the case in which the
residual-dependent term of Proposition~\ref{prop:confound} vanishes. By part
(a), neither control depends on the value of the constant it holds fixed,
except through $\varepsilon_{\mathrm{A}}$ and clipping; only the variation of
the weights across timesteps and bands matters.

For AWWL at $\alpha = 0.2$, $p = 1$ on the CIFAR-10 schedule, holding $S$
fixed leaves $\sum_b \omega_b^2$ varying by a factor of $4.00 = J$ over the
training timesteps. Both bounds of Eq.~\eqref{eq:cs-bounds} are approached:
the left one at $\lambda_p(\sigma) = (1-\alpha)/(1+2\alpha)$, that is
$\sigma = 4/7$, where the four per-band weights are equal, and the right one as
$\sigma \to 1$, where all weight moves to the approximation band.

\subsection{The residual spectrum under a Gaussian model}
\label{app:residual}

\begin{proposition}[Residual of the optimal predictor]
\label{prop:residual}
Let $x_0 \sim \mathcal{N}(0, C)$ with $C = U \operatorname{diag}(P)\, U^\top$ for
an orthonormal basis $U$. For a stationary (circulant) covariance, $U$ can be
taken to be the real Fourier basis and $P$ is the power spectrum. Let
$\hat\epsilon = \mathbb{E}[\epsilon \mid x_t]$ be the optimal predictor at
timestep $t$ (Proposition~\ref{prop:minimiser}). Then the coordinates of
$U^\top(\hat\epsilon - \epsilon)$ are independent, with variances
\begin{equation}
  v_t(u) = \frac{\bar\alpha_t P(u)}{\bar\alpha_t P(u) + \sigma_t^2} .
  \label{eq:gaussian-residual}
\end{equation}
For two coordinates with $P(u) > P(u') > 0$, the ratio $v_t(u)/v_t(u')$
increases strictly from $1$ to $P(u)/P(u')$ as $\mathrm{SNR}_t$ decreases from
$\infty$ to $0$.
\end{proposition}

\begin{proof}
Write $z = U^\top x_0$, $\tilde\epsilon = U^\top \epsilon$ and
$y = U^\top x_t = \sqrt{\bar\alpha_t}\, z + \sigma_t \tilde\epsilon$. Because $U$ is
orthonormal, $\tilde\epsilon \sim \mathcal{N}(0, I)$, and $z$ and $\tilde\epsilon$ are independent with
independent coordinates, so the problem decouples across coordinates. For one
coordinate, $(\tilde\epsilon_u, y_u)$ is jointly Gaussian with
$\operatorname{Var}(y_u) = \bar\alpha_t P(u) + \sigma_t^2$ and
$\operatorname{Cov}(\tilde\epsilon_u, y_u) = \sigma_t$, so
$\mathbb{E}[\tilde\epsilon_u \mid y] = \sigma_t\, y_u/(\bar\alpha_t P(u) + \sigma_t^2)$ and
the residual variance is
$1 - \sigma_t^2/(\bar\alpha_t P(u) + \sigma_t^2) = v_t(u)$; each residual
coordinate depends only on $(z_u, \tilde\epsilon_u)$, hence independence. With
$s = 1/\mathrm{SNR}_t = \sigma_t^2/\bar\alpha_t$,
$v_t(u)/v_t(u') = \big(P(u)/P(u')\big)\,\big(P(u') + s\big)/\big(P(u) + s\big)$,
whose derivative in $s$ has the sign of $P(u) - P(u') > 0$; it equals $1$ at
$s = 0$ and tends to $P(u)/P(u')$ as $s \to \infty$.
\end{proof}

For natural images $P$ decreases with frequency, so the optimal predictor's
residual is nearly white at low noise ($v_t \to 1$) and shifts its energy
towards low frequencies as the noise level grows. The normalised residual
spectra $v_t/\sum_u v_t(u)$ are ordered in the likelihood-ratio sense in $P$
(for $\sigma_{t'} > \sigma_t$ the ratio $v_{t'}(u)/v_t(u)$ increases with
$P(u)$), so the share of residual energy in any band whose squared frequency
response is an increasing function of $P$ grows with $\sigma_t$. The Haar
approximation band is not exactly of this form, because its response is not a
function of radial frequency alone, but its response is concentrated at the
low frequencies where a natural-image spectrum is largest; evaluating
$\bar\rho_{\mathrm{LL}}(t)$ directly for $P(u) \propto \lvert u \rvert^{-2}$ on
the CIFAR-10 schedule gives a monotone increase from $1$ at the lowest noise
level. In Proposition~\ref{prop:confound} this makes $\Lambda(t)$ depend on the
timestep even when $S$ is held fixed. How far a trained model's residual
follows this pattern is an empirical question.

\subsection{Gradient-norm bounds}
\label{app:gradnorm}

The two controls of Section~\ref{sec:confound-controls} hold different
quantities fixed because the loss scale and the gradient scale of a band
weighting differ. For one sample,
$\nabla_\theta \ell_\sigma = \sum_b \omega_b(\sigma)\, q_b$, with $q_b$ the
parameter gradient of the mean square of band $b$. If the $q_b$ have a common
norm $G$ and non-negative inner products, Proposition~\ref{prop:scale}(b)
gives
\begin{equation}
  G \Big(\sum_b \omega_b(\sigma)^2\Big)^{1/2}
  \;\le\; \lVert \nabla_\theta \ell_\sigma \rVert
  \;\le\; G\, S(\sigma).
  \label{eq:gradscale}
\end{equation}
The upper bound holds with equality when the band gradients are parallel, the
lower when they are orthogonal, as the output-space band gradients are
(Proposition~\ref{prop:scale}(c)). Holding $S$ fixed therefore fixes the
per-timestep gradient weight in the parallel case and, for residuals with
equal band energies, the output-space component along the pixel-MSE gradient
(Eq.~\eqref{eq:mse-component}); holding $\sum_b \omega_b^2$ fixed fixes it in
the orthogonal case. By Eq.~\eqref{eq:cs-bounds} with $J = 4$ bands, fixing
either quantity leaves the other free within a factor of $\sqrt{J} = 2$ in
norm, and both limits are approached on the CIFAR-10 schedule
(Fig.~\ref{fig:gradient-scale}; the points at which they are reached are given
at the end of Appendix~\ref{app:scale}). Neither control is therefore the
stronger: each fixes one end of Eq.~\eqref{eq:gradscale}. The first control
multiplies the AWWL loss by $c/S(\sigma)$ and the second by
$\big(\levelSqFixed / \sum_b \omega_b(\sigma)^2\big)^{1/2}$. Both statements
about the gradient weight assume $\rho_b(r) = 1$; for other residuals the
residual-dependent term of Proposition~\ref{prop:confound} remains, and
neither control fixes the effective timestep weight $\Lambda(t)$.

\guardedfigure{1.0}{gradient_scale.png}{%
  What each control holds fixed, for the released form of AWWL at
  $\alpha = 0.2$, $p = 1$, with per-band weights
  $\omega_b = (w_{\mathrm{LL}}, w_{\det}/3, w_{\det}/3, w_{\det}/3)$ as in
  the panel titles. (a) The total
  $S = \sum_b \omega_b$, which sets the loss scale at each timestep and the
  gradient weight when the band gradients are parallel. (b)
  $\sum_b \omega_b^2$, which sets the squared gradient weight when the band
  gradients are orthogonal with equal norms (Eq.~\eqref{eq:gradscale}). Each
  curve is divided by AWWL's mean of the same quantity over the training
  timesteps (\wTotalMean{} in (a), \wSqMean{} in (b)). Solid: AWWL; dotted:
  total fixed at \levelSumFixed; dash-dotted: sum of squares fixed at
  \levelSqFixed. Fixing either quantity leaves the other varying, by a factor
  of $4$ in (b) (of $2$ in its square root) and of $2$ in (a). Horizontal
  axes: training timestep $t$, sampled uniformly; top axes:
  $\sigma_t$.}{fig:gradient-scale}
}{%
  \todo{sections/appendix\_theory.tex has not been written yet.}}

\FloatBarrier
\section{Experimental details}\label{app:setup}
\IfFileExists{sections/appendix_setup.tex}{% Appendix B, experimental details. appendix.tex carries \section and
% \label{app:setup}, so this file starts at subsection level. Every value below
% is a design constant read from the released configuration files (configs/)
% or from the code defaults they rely on; measured quantities come from
% generated/numbers.tex.

\subsection{CIFAR-10 recipe}
\label{app:setup-cifar}

Table~\ref{tab:setup} lists the settings of the CIFAR-10 matrix and of the
settings derived from it. They are those of the released configuration files,
and every run stores its resolved configuration next to its checkpoints.

\begin{table}[tbp]
\centering
\footnotesize
\setlength{\tabcolsep}{3pt}
\caption{Training and evaluation settings. ``$=$'' means as in the CIFAR-10
matrix. $s$ is the spectral statistic of Section~\ref{sec:pricing-test}. The
learning rates of objectives with learnable parameters are given in
Table~\ref{tab:arms}; the DreamBooth recipe is given in
Appendix~\ref{app:setup-dreambooth}.}
\label{tab:setup}
\begin{tabular}{@{}>{\raggedright\arraybackslash}p{2.2cm}
                >{\raggedright\arraybackslash}p{4.4cm}
                >{\raggedright\arraybackslash}p{2.75cm}
                >{\raggedright\arraybackslash}p{2.75cm}
                >{\raggedright\arraybackslash}p{2.75cm}@{}}
\toprule
 & CIFAR-10 matrix & Four times the budget & CelebA-64 & Denoising \\
\midrule
\multicolumn{5}{@{}l}{\emph{Data and model}} \\
Training set & CIFAR-10 train \citep{krizhevsky2009learning} & $=$
  & CelebA train \citep{liu2015deep}, \num{162770} images & $=$ \\
Resolution & $32\times32$ & $=$ & $64\times64$ (short side resized,
  centre crop) & $=$ \\
Preprocessing & scaled to $[-1,1]$, random horizontal flip & $=$ & $=$ & $=$ \\
Network & U-Net, channels (128, 256, 256, 256), two residual blocks per
  level, one self-attention layer (bottleneck only), dropout~0, 34.4\,M
  parameters & $=$ & $=$ & $=$; maps the degraded image to the clean one \\
Task & $T = 1000$, $\beta_t$ linear from $10^{-4}$ to $0.02$,
  $\epsilon$-prediction, $t$ uniform & $=$ & $=$ & additive Gaussian noise,
  $\sigma \sim \mathcal{U}[0.05, 0.5]$ on the $[-1,1]$ scale; $\sigma$ drives
  the time embedding \\
\midrule
\multicolumn{5}{@{}l}{\emph{Optimisation}} \\
Optimiser & AdamW \citep{loshchilov2019decoupled}, learning rate
  $2\times10^{-4}$, $(\beta_1,\beta_2) = (0.9, 0.999)$, $\varepsilon_{\mathrm{A}} = 10^{-8}$,
  weight decay $0.01$ & $=$, fused kernel & $=$ & $=$ \\
Schedule & 500 linear warm-up steps, then cosine decay to zero & $=$, over
  800 epochs & $=$ & $=$ \\
Batch; epochs; steps & 128; 200; \num{78200} & 128; 800; \num{312800}
  & 128; 200; \num{254200} & 128; 100; \num{39100} \\
Clipping; precision & gradient norm 1.0; float16 mixed precision & $=$ & $=$
  & $=$ \\
Weight averaging & none (EMA arms: decay 0.9999) & EMA, decay 0.9999 & none
  & none \\
\texttt{torch.compile} & off & on & off & off \\
Arms; seeds & Table~\ref{tab:arms}; 1--5 (1--3 for four arms) & MSE, AWWL;
  1--3 & MSE; 1--5 & six objectives; 1--3 \\
\midrule
\multicolumn{5}{@{}l}{\emph{Evaluation}} \\
Checkpoints & epoch 199; also 39, 79, 119, 159 for MSE, static, AWWL
  & epochs 199 to 799, every 100 & epoch 199 & epoch 99 \\
Sampler & DDIM-100, $\eta = 0$; DDPM-1000 at epoch 199 for MSE, static, AWWL
  & DDIM-100; DDPM-1000 at epoch 799 & DDIM-100 & one forward pass \\
Samples; reference & \num{10000}; all \num{\nRefImages} training images & $=$
  & \num{10000}; all \num{162770} training images & 500 test images \\
Metrics & FID, IS, KID, precision, recall; $s$ & FID, IS, KID; $s$ & FID;
  $s$ & PSNR, SSIM, LPIPS, NIQE \\
\bottomrule
\end{tabular}
\end{table}

\paragraph{Network.}
The denoiser is \texttt{UNet2DModel} of \texttt{diffusers} with plain residual
blocks (\texttt{DownBlock2D} and \texttt{UpBlock2D}) at every resolution level
and the library's default bottleneck, two residual blocks around one
self-attention layer with 32 heads of dimension 8. It uses group normalisation
with 32 groups, SiLU activations and a sinusoidal timestep embedding projected
to 512 dimensions. The first three levels halve the resolution with strided
convolutions, so the bottleneck operates at $4\times4$ on CIFAR-10 and at
$8\times8$ on CelebA-64; the parameter count does not depend on the
resolution.

\paragraph{Sampling.}
DDIM uses 100 timesteps spaced by 10 from $t = 990$ down to $t = 0$, with
$\eta = 0$ and the predicted clean image clipped to $[-1,1]$. The ancestral
sampler uses all 1000 steps with the posterior variance of the forward
process. Samples are generated in batches of 256 from the sampling seed 12345,
the same for every arm, training seed and checkpoint. The assignment of noise
to images depends on the batch size, so the batch size is the same in every
sampling run. Arms trained with EMA are sampled from the averaged weights.

\paragraph{Execution.}
Compilation with \texttt{torch.compile} and the fused AdamW kernel change
kernel selection and the order of floating-point operations, so a compiled arm
would not be exactly comparable with an uncompiled baseline. Both are off
throughout the 200-epoch matrix. The four-times-budget runs use both, because
both of their arms were trained for that comparison and neither is compared
with an uncompiled run of the same budget.

\subsection{Arms of the CIFAR-10 matrix}
\label{app:setup-arms}

Table~\ref{tab:arms} lists the arms with their settings. The timestep-only
arm separates AWWL's timestep component from its frequency shape
(Section~\ref{sec:confound}); it belongs to neither the main nor the pricing
family but to a two-test family of its own (Appendix~\ref{app:setup-stats}). Every wavelet arm except the learned basis
and the combination of all axes, which use the lifting transform, uses the
Haar DWT with zero boundary
extension. Arms with learnable loss parameters optimise them with AdamW in a
parameter group of their own, at the learning rate given in the table and
under the same schedule as the network.

The VGG perceptual loss of the conference paper is not replicated. The
trainer applies every objective to the pair $(\hat\epsilon, \epsilon)$, and
the target $\epsilon$ is white Gaussian noise, whose VGG features carry no
perceptual content, so the loss would not measure what its name suggests; at
the $224\times224$ input of VGG-16 it would also upsample every $32\times32$
batch 49-fold, twice per step.

\begin{table}[tbp]
\centering
\footnotesize
\setlength{\tabcolsep}{4pt}
\caption{Arms of the CIFAR-10 matrix, grouped and named as in the result
tables; MSE is the baseline. $\omega_b$ are the per-band weights of
Eq.~\eqref{eq:family}; ``AWWL weighting'' is the averaged-detail weighting of
Eq.~\eqref{eq:audited} and ``summed weighting'' the summed-detail weighting of
Eq.~\eqref{eq:summed}, both with $\alpha = 0.2$ and $p = 1$ unless stated. A
constraint on $\sum_b \omega_b$ or $\sum_b \omega_b^2$ holds at every noise
level and rescales the weights without changing their ratios; the four
normalisation controls impose such a constraint. Learning rates are those of
the learnable loss parameters.}
\label{tab:arms}
\begin{tabular}{@{}>{\raggedright\arraybackslash}p{4.6cm}
                >{\raggedright\arraybackslash}p{8.9cm}r@{}}
\toprule
Arm & Objective and settings & Seeds \\
\midrule
MSE & $\mathcal{L}_{\mathrm{simple}}$, Eq.~\eqref{eq:simple} & 5 \\
\addlinespace[2pt]
\multicolumn{3}{@{}l}{\emph{Conference-paper arms}} \\
AWWL & AWWL weighting, one-level Haar DWT & 5 \\
Static wavelet ($p = 0$) & AWWL weighting with $p = 0$ & 5 \\
AWWL, detail bands summed & summed weighting & 5 \\
\addlinespace[2pt]
\multicolumn{3}{@{}l}{\emph{Normalisation controls}} \\
AWWL, $\sum_b \omega_b = 1$ & AWWL weighting, $\sum_b \omega_b = 1$ & 5 \\
AWWL, $\sum_b \omega_b = 0.5$ & AWWL weighting, $\sum_b \omega_b = 0.5$ & 5 \\
AWWL, $\sum_b \omega_b^2$ fixed & AWWL weighting,
  $\sum_b \omega_b^2 = \wSqMean$, the mean of $\sum_b \omega_b^2$ under the
  AWWL weighting over the training timesteps & 5 \\
AWWL, detail bands summed, $\sum_b \omega_b^2$ fixed & summed weighting,
  $\sum_b \omega_b^2 = \wSqMeanSummed$, the mean of $\sum_b \omega_b^2$ under
  the summed weighting over the training timesteps & 5 \\
\addlinespace[2pt]
\multicolumn{3}{@{}l}{\emph{Timestep schedule alone}} \\
AWWL, timestep only & $S(\sigma_t)\,\mathcal{L}_{\mathrm{simple}}$, with $S$
  the total weight of AWWL (Eq.~\eqref{eq:total}) and no frequency shape & 5 \\
\addlinespace[2pt]
\multicolumn{3}{@{}l}{\emph{Planted-mechanism control}} \\
MSE, planted change & $\mathcal{L}_{\mathrm{simple}}$ on training images whose
  spectrum carries AWWL's measured change: each channel's centred DFT is
  multiplied by $10^{\bar\delta(r)/20}$ at integer radius $r$, with
  $\bar\delta$ the seed mean of $s(A) - s(B)$ at epoch 199, in floating point
  and clipped to $[0,1]$ before normalisation & 5 \\
\addlinespace[2pt]
\multicolumn{3}{@{}l}{\emph{Other weightings}} \\
Min-SNR-$\gamma$ & $\kappa_t = \min(\mathrm{SNR}_t, \gamma)/\mathrm{SNR}_t$,
  $\gamma = 5$ \citep{hang2023efficient} & 5 \\
Raw SNR weighting & timestep weight $\kappa_t = \mathrm{SNR}_t$
  (Section~\ref{sec:notation}) & 5 \\
FFL & Eq.~\eqref{eq:ffl} with $a = 1$, weights computed per image
  \citep{jiang2021focal} & 5 \\
\addlinespace[2pt]
\multicolumn{3}{@{}l}{\emph{Pixel losses}} \\
$L_1$ & mean absolute residual & 5 \\
Huber & Huber loss with threshold 1 \citep{huber1964robust} & 5 \\
Charbonnier & mean of $\sqrt{r^2 + \varepsilon_{\mathrm{C}}^2}$, $\varepsilon_{\mathrm{C}} = 10^{-6}$
  \citep{charbonnier1997deterministic} & 5 \\
\addlinespace[2pt]
\multicolumn{3}{@{}l}{\emph{Extensions (unnormalised)}} \\
Per-sub-band schedules & summed weighting, $p = 2$ for $\mathrm{HH}$ and
  $p = 1$ for $\mathrm{LH}$ and $\mathrm{HL}$ & 5 \\
Spatially adaptive & summed weighting; each detail weight is multiplied by a
  map of local activity of the target's coefficients ($3\times3$ mean of their
  magnitude), raised to the power $1-\sigma_t$ and renormalised to mean 1 per
  image and band; the target is the noise $\epsilon$, so the map carries no
  image structure & 5 \\
Multi-level (3 levels) & three-level Haar DWT; the approximation band keeps
  $w_{\mathrm{LL}}$ and $w_{\det}$ is divided equally among the nine detail
  bands & 5 \\
Learned weighting & uncertainty weighting \citep{kendall2018multi}: a small
  network maps $\sigma_t$ to one log-variance per band; learning rate
  $10^{-3}$ & 5 \\
GradNorm balancing & band weights set by GradNorm \citep{chen2018gradnorm}
  on the output layer, asymmetry 1.5, weights renormalised to a fixed sum
  after every step; learning rate $2.5\times10^{-2}$ & 5 \\
Learned basis (lifting) & summed weighting on a learnable lifting transform
  \citep{sweldens1998lifting} initialised at the Haar wavelet; its filters are
  trained on the loss they define, with no orthogonality constraint; learning
  rate $10^{-4}$ & 3 \\
All axes combined & learned weighting, spatial modulation and lifting basis
  together; learning rate $10^{-3}$ & 3 \\
\addlinespace[2pt]
\multicolumn{3}{@{}l}{\emph{With EMA}} \\
MSE + EMA & MSE; EMA of the weights, decay 0.9999 & 3 \\
AWWL + EMA & AWWL; EMA of the weights, decay 0.9999 & 3 \\
\bottomrule
\end{tabular}
\end{table}

\subsection{Four times the training budget}
\label{app:setup-budget}

The four-times-budget runs train MSE and AWWL for 800 epochs with EMA, three
seeds each, with the settings of the corresponding column of
Table~\ref{tab:setup}. The cosine
schedule spans the whole run, so a longer run is also annealed more slowly: at
epoch 199 these runs are at about \SI{85}{\percent} of the peak learning rate,
whereas the 200-epoch runs have reached zero. Their intermediate checkpoints
are therefore points of a different schedule, and a comparison across budgets
is between endpoints: epoch 799 here against epoch 199 of the EMA arms of the
matrix. Such a comparison also spans the change of execution path described
above, and is descriptive; the tests compare AWWL with MSE within one budget.
The same weights at epoch 799 are also sampled with DDPM-1000.

\subsection{CelebA-64}
\label{app:setup-celeba}

The CelebA-64 runs keep the network and the optimisation of the CIFAR-10
matrix, including the 200 epochs, so they change more than the data and the
resolution. The larger training set gives \num{254200} optimiser steps
against \num{78200}, 3.25 times as many, and the higher resolution also
enlarges the bottleneck to $8\times8$ and lowers the factor by which FID
upsamples the images to $299\times299$ from 9.3 to 4.7. The training set is
the official CelebA training partition; each image is resized so that its
short side is 64 pixels and centre-cropped to $64\times64$, which preserves
the aspect ratio, and the FID reference set is the same partition prepared in
the same way. The last incomplete batch of each epoch is dropped, which gives
1271 steps per epoch; the CIFAR-10 loader keeps it (391 steps per epoch).
Only MSE is trained, with five seeds; the pricing test on this dataset
applies its operators to MSE's samples (Section~\ref{sec:celeba}). The
spectral statistic uses the same definition on the $64\times64$ grid, with
the bands again the thirds of the radial bins.

\subsection{DreamBooth}
\label{app:setup-dreambooth}

DreamBooth \citep{ruiz2023dreambooth} fine-tunes the full U-Net of Stable
Diffusion~1.5 \citep{rombach2022high} on one subject of the DreamBooth dataset
(\emph{robot toy}, five instance images) with the instance prompt ``a photo of
sks robot toy''. The images are resized and centre-cropped to $512\times512$.
The VAE and the text encoder are frozen, and no prior-preservation loss is
used. Training runs for 400 steps at batch size 1 with AdamW at a constant
learning rate of $5\times10^{-6}$, gradient-norm clipping at 1.0 and float16
mixed precision, under the model's own noise schedule ($T = 1000$,
scaled-linear $\beta_t$ from $0.00085$ to $0.012$, $\epsilon$-prediction). The
objectives therefore act on the residual of the four-channel $64\times64$
latent. Ten objectives are compared: MSE, $L_1$, Huber, Charbonnier, raw SNR
weighting, Min-SNR-$\gamma$ ($\gamma = 5$), the static wavelet arm ($\alpha = 0.8$,
$p = 0$), AWWL with its published DreamBooth setting ($\alpha = 0.8$, $p = 2$),
AWWL normalised to $\sum_b \omega_b = 1$, and the learned weighting (learning
rate $10^{-3}$), each with five seeds.

Each run generates 20 images for each of three prompts (the subject in a
cyberpunk neon city, on a beach at sunset, and as a watercolour painting) with
the pipeline's default PNDM scheduler, 50 steps, classifier-free guidance
scale 7.5 and image seeds 1 to 20. CLIP score is the cosine similarity between
the CLIP ViT-B/32 \citep{radford2021learning} embeddings of the prompt and the
image; subject similarity is the cosine similarity between the embedding of
the image and the mean embedding of the instance images. Both are averaged
over the 60 images of a run. Radial spectra of generated and instance images
are computed after resizing both to $512\times512$, because DFT magnitudes
scale with the number of pixels and profiles taken at different resolutions
are not comparable.

\subsection{Gaussian denoising}
\label{app:setup-restoration}

The denoising task trains the CIFAR-10 U-Net as a one-step restorer: the input
is a training image corrupted by additive Gaussian noise of standard
deviation $\sigma \sim \mathcal{U}[0.05, 0.5]$ on the $[-1,1]$ intensity
scale, clipped to that range, and the target is the clean image. The same
$\sigma$ drives the band weights of the wavelet objectives and, mapped
linearly onto the timestep range, the network's time embedding. Six objectives
are compared with three seeds each: MSE, $L_1$, Charbonnier, the static
wavelet arm ($\alpha = 0.5$, $p = 0$), AWWL at the balanced setting $\alpha = 0.5$,
$p = 1$, and AWWL normalised to $\sum_b \omega_b = 1$. Evaluation uses 500
images drawn from the CIFAR-10 test set, with noise levels evenly spaced over
$[0.05, 0.5]$ and assigned in random order; the selection and the assignment
are fixed by one seed and identical for every arm. PSNR uses peak value 1
after rescaling to $[0,1]$; SSIM
\citep{wang2004image} uses an $11\times11$ Gaussian window with standard
deviation 1.5; LPIPS \citep{zhang2018unreasonable} uses the AlexNet backbone.
NIQE \citep{mittal2013making} is implemented from its definition, with the
pristine model fitted to the 500 clean evaluation images and the statistics of
the restored set pooled over all its images, so its values are not comparable
with NIQE values computed against the standard pristine corpus.

\subsection{Metric implementations and software}
\label{app:setup-metrics}

FID uses \texttt{clean-fid} in its \texttt{clean} mode, with its port of the
Inception-V3 network \citep{szegedy2016rethinking} and resizing to
$299\times299$ inside the library. Inception Score (ten splits) and KID (100
subsets of \num{1000} images, cubic polynomial kernel, against all reference
images) come from \texttt{torch-fidelity} \citep{obukhov2020torchfidelity}.
Precision and recall \citep{kynkaanniemi2019improved} come from
\texttt{prdc} \citep{naeem2020reliable} with $k = 5$ nearest neighbours, on the
2048-dimensional pool features of the ImageNet-trained Inception-V3 of
\texttt{torchvision}, comparing the \num{10000} samples with the first
\num{10000} training images in sorted file order. The wavelet transforms use
\texttt{pytorch-wavelets}. The reported numbers were computed with PyTorch
2.13, \texttt{diffusers} 0.39, \texttt{clean-fid} 0.1.35,
\texttt{torch-fidelity} 0.4.0, \texttt{prdc} 0.2, \texttt{pytorch-wavelets}
1.3.0 and \texttt{lpips} 0.1.4.

\subsection{Statistical analysis}
\label{app:setup-stats}

Table~\ref{tab:families} lists the Holm families. Each family answers one
question, and each hypothesis belongs to exactly one family. The budget family
holds only the two 800-epoch contrasts; the 200-epoch operating points shown
beside them in Table~\ref{tab:budget} are tested in the main, EMA and sampler
families. The decomposition family tests the timestep-only arm against MSE,
the effect of AWWL's timestep component, and AWWL against that arm, the effect of adding the frequency shape to the same
total weight $S(\sigma_t)$. The arm's MRG lies outside every family and is
reported with its raw $p$ only (Table~\ref{tab:pricing-arms}).

\begin{table}[tbp]
\centering
\small
\caption{Holm families. Each family, and within it each metric, is corrected
separately; each hypothesis belongs to one family.}
\label{tab:families}
\begin{tabular}{@{}l>{\raggedright\arraybackslash}p{0.55\linewidth}rl@{}}
\toprule
Family & Hypotheses & Size & Reported in \\
\midrule
Main & each arm against MSE, CIFAR-10, epoch 199, DDIM-100 (with MSE + EMA;
  without AWWL + EMA and the timestep-only arm) & \nFamilyMain
  & Table~\ref{tab:compare-fid} \\
EMA & AWWL + EMA against MSE + EMA & 1 & Table~\ref{tab:compare-fid} \\
Confound & each normalisation control against the unnormalised AWWL with the
  same detail reduction & \nFamilyConfound & Table~\ref{tab:confound} \\
Decomposition & timestep-only arm against MSE; AWWL against the timestep-only
  arm & 2 & Table~\ref{tab:confound} \\
Extensions & each extension against AWWL with detail bands summed
  & \nFamilyExtensions & Table~\ref{tab:extensions} \\
Sampler & static wavelet arm and AWWL against MSE, DDPM-1000, epoch 199 & 2
  & Table~\ref{tab:sampler} \\
Budget & AWWL against MSE after 800 epochs, DDIM-100 and DDPM-1000 & 2
  & Table~\ref{tab:budget} \\
Pricing & MRG $= 0$ for the transplant of each stably trained arm of the
  main family and of AWWL + EMA (Section~\ref{sec:worth-arms}) & \nFamilyPricing
  & Table~\ref{tab:pricing-arms} \\
DreamBooth & each objective against MSE & 9 & Table~\ref{tab:dreambooth} \\
Denoising & each objective against MSE & 5 & Table~\ref{tab:restoration} \\
\bottomrule
\end{tabular}
\end{table}

Pairing follows from the initialisation and data order that seed $k$ fixes
for every arm (the sampling latents are common to all runs), and it reduces
variance only if two arms' per-seed outcomes correlate positively. The
estimated correlation between the per-seed FID of AWWL and of MSE is
$\rSeedAWWL$ under DDIM-100 and $\rSeedAWWLddpm$ under DDPM-1000; with five pairs each, both are compatible with zero. The paired
test is exact whatever the correlation. Near zero correlation it has fewer
degrees of freedom than an unpaired test, and under a negative correlation an
unpaired test would understate the variance of the difference and so be
anti-conservative. The paired test is therefore used throughout.

A non-significant difference is read through the effects its interval
excludes, not as evidence of no effect. At five seeds and the observed
standard deviation of the paired AWWL$-$MSE differences
(\sdPairedAWWL~FID), a two-sided paired test at the $0.05$ level has 80\%
power against a true difference of \mdeAWWL~FID. This describes the resolution
of the design rather than a prior power calculation, and it is itself
imprecise, because the standard deviation is estimated on four degrees of
freedom.

\subsection{Compute}
\label{app:setup-compute}

Training used NVIDIA RTX 5000 Ada GPUs (32~GB), one run per GPU. The
\gpuHoursTotal{} GPU-hours of training over \nRunsTotal{} runs divide into
\gpuHoursCifar{} for the \nRunsCifar{} runs of the CIFAR-10 matrix,
\gpuHoursLong{} for the \nRunsLong{} four-times-budget runs,
\gpuHoursCeleba{} for the \nRunsCeleba{} CelebA-64 runs,
\gpuHoursRestoration{} for the \nRunsRestoration{} denoising runs and
\gpuHoursDreambooth{} for the \nRunsDreambooth{} DreamBooth runs. Sampling,
scoring and the pricing passes are not included in these figures.
}{%
  \todo{sections/appendix\_setup.tex has not been written yet.}}

% appendix_scoring.tex has no heading of its own; its \label{app:scoring}
% takes the number of this subsection.
\subsection{Implementation of the pricing test}
\IfFileExists{sections/appendix_scoring.tex}{% Implementation of the pricing test: the statistic, the boost, the scoring
% path and the status of the preconditions (Table tab:preconditions), moved
% from or cited by Section~\ref{sec:pricing-test}. No heading: the integrator
% inputs this file inside Appendix~B (Experimental details), e.g. at the end
% of the subsection on metric implementations. \label{app:scoring} takes the
% number of the enclosing (sub)section and is cited from pricing.tex.
\label{app:scoring}%
The statistic of the pricing test, Eq.~\eqref{eq:statistic}, is computed on
the luminance
$y = 0.299\,x_{\mathrm{R}} + 0.587\,x_{\mathrm{G}} + 0.114\,x_{\mathrm{B}}$ of
an RGB image $x$ with values in $[0, 255]$ (ITU-R BT.601 weights, unrounded).
Its 2-D DFT $\hat y = F y$ on the $h \times w$ pixel grid is indexed by $u$ after the
zero frequency is shifted to $c = (\lfloor h/2 \rfloor, \lfloor w/2 \rfloor)$;
the integer radius is $r(u) = \lfloor \lVert u - c \rVert_2 \rfloor$, and the
bins are $U_k = \{u : r(u) = k\}$, $k = 0, \dots, K-1$. Averaging decibels
rather than power makes $s$ a geometric-mean spectrum, and the normalisation of
$F$ cancels in every difference used. Because $s$ sees only the luminance, the
operators impose its change on every colour channel, and a change in chroma is
neither measured nor imposed separately. The boost is the gain
$g = 10^{d \varrho / 20}$ of dose $d$ in decibels, with the raised-cosine ramp
$\varrho = \frac12 - \frac12 \cos\big(\pi \min\{1, \max\{0, (\nu - 0.667)/0.5\}\}\big)$
in the radius $\nu = \lVert u - c \rVert_2 / (\min(h, w)/2)$, measured in units
of the axial Nyquist frequency. White noise $\xi$ is drawn independently for
every pixel and channel and added without rounding.

Every term of Eqs.~\eqref{eq:value}--\eqref{eq:accounting} is scored in
floating point. The filtered images, clipped to $[0, 255]$ and not rounded,
enter the `clean' resizer and feature extractor of \texttt{clean-fid}
\citep{parmar2022aliased} directly: bicubic resizing of each channel to
$299 \times 299$ in floating point, 2048-dimensional Inception-V3 pool
features \citep{szegedy2016rethinking}, and the Fr\'echet distance to the
reference features \citep{heusel2017gans}. FD$_{\text{DINOv2}}$
(Section~\ref{sec:worth-dino}) replaces the extractor by DINOv2 ViT-L/14
\citep{oquab2024dinov2}: images are resized bicubically to $224 \times 224$ in
floating point, normalised with the ImageNet mean and standard deviation, and
represented by the model's normalised class token, as in
\citet{stein2023exposing} except that the resize acts on floating-point rather
than 8-bit images, for original and filtered samples alike. For unmodified samples the path
reproduces \texttt{clean-fid}'s file-based computation: the baseline FIDs
agree with those of the training ledger to within \priceIdentityMaxDiff. The
identity operator used for this check bypasses the filter and the clipping,
so the filter path at zero change is exercised only by arms whose change is
close to zero, such as the static weighting of the null control. Runs that
reproduce an earlier 8-bit measurement round the filtered images, with or
without a dither of $\pm 0.5$~LSB, to isolate the share of the value due to
re-quantisation (hazard H1). The test costs one computation of $s$ per set
and one feature-extraction pass per operator and seed.

Two effects of the 8-bit input remain in the floating-point values. The
operators act on $B$'s stored 8-bit samples, so a gain $g$ also rescales $B$'s
own rounding-noise field; in each bin this is equivalent, in power, to
adding independent noise of $g^2 - 1$ times the rounding variance, a field of
the kind the white-noise control prices. Its share of $V$ is not separated, and
it grows with $g^2 - 1$, most for the larger boost doses and the match.
Clipping, in turn, makes the change measured back after an operator differ
slightly from the imposed $\delta_i$, most for arms whose change is large;
Table~\ref{tab:pricing-arms} reports the change measured back.

Table~\ref{tab:preconditions} lists where each precondition of
Section~\ref{sec:pricing-pre} is checked and whether it holds in the three
settings.

\begin{table}[tbp]
\centering
\small
\caption{Status of the preconditions of the pricing test. Section: where the
  CIFAR-10 evidence is; that for CelebA-64 is in Section~\ref{sec:celeba}, that
  for DreamBooth in Section~\ref{sec:dreambooth}. P3: on CIFAR-10 every
  spectral operator is a coherent radial gain and only the transplant imposes
  $\delta_A$; on CelebA-64 the coherent gains and white noise move FID in
  opposite directions. n/a: needs the AWWL
  arm, which was not trained on CelebA-64; not run: the check was not
  performed.}
\label{tab:preconditions}
\begin{tabular}{@{}llllll@{}}
\toprule
 & Requirement & Section & CIFAR-10 & CelebA-64 & DreamBooth \\
\midrule
P0 & no re-quantisation & \ref{sec:worth-curve} & met & met
  & not met (8-bit) \\
P1 & $s$ moves as scored & \ref{sec:mechanism} & DDIM-100 only & n/a
  & met \\
P2 & $\delta$ measured back & \ref{sec:worth-headline} & met (transplant)
  & n/a & not run \\
P3 & several operators & \ref{sec:worth-checks}
  & coherent gains only & signs differ & boost only \\
P4 & metric registers $s$ & \ref{sec:worth-curve}
  & met & met & not met \\
P5 & no interaction & \ref{sec:worth-checks} & consistent & n/a
  & not run \\
\bottomrule
\end{tabular}
\end{table}
}{%
  \todo{sections/appendix\_scoring.tex has not been written yet.}}

\FloatBarrier
\section{Pricing and comparison tables}\label{app:tables}
This appendix collects the generated tables behind the pricing test and the
five-seed comparison. Each is generated from the run ledgers by the scripts in
the repository, and its caption states how its entries are computed,
including the test and the correction family behind any $p$-value.

Table~\ref{tab:price-curve} gives the value of the unquantised high-band
boost at each dose on MSE's CIFAR-10 samples (Section~\ref{sec:worth-curve})
and the prices of the boost and of the radial match on CelebA-64
(Section~\ref{sec:celeba}). Table~\ref{tab:operators} lists each operator
applied to MSE's samples with its value and the high-band deficit it leaves,
including the transplant of AWWL's change and the 8-bit variants behind
hazard H1 (Section~\ref{sec:worth}). Table~\ref{tab:pricing-arms} gives the
high-band change, gap, value and MRG of each of the \nFamilyPricing{} arms of
the pricing family, priced against MSE, which Fig.~\ref{fig:pricing-arms}
plots (Section~\ref{sec:worth-arms}). Table~\ref{tab:worth} applies the
$+0.41$~dB high-band boost to the samples of MSE and of AWWL, the secondary
check of precondition P5; the main check, the reverse transplant, is
reported in Section~\ref{sec:worth-checks}.

Table~\ref{tab:features} repeats the pricing test in DINOv2 features
(Section~\ref{sec:worth-dino}) beside the same quantities in FID.
Table~\ref{tab:compare-fid} gives the paired comparison with MSE on FID that
Fig.~\ref{fig:effects} plots (Section~\ref{sec:compare}), with the raw,
Holm-corrected and Wilcoxon $p$-values of every arm, and
Table~\ref{tab:compare-other} repeats it on Inception Score, KID, precision
and recall, with the Holm correction applied within each metric.
On Inception Score AWWL does not differ from MSE, and every objective that
differs after correction scores lower. On KID neither AWWL nor
Min-SNR-$\gamma$ nor FFL differs from MSE, although the latter two are worse
on FID. Precision is the only metric on which any objective improves on MSE
after correction (Min-SNR-$\gamma$ and the learned weighting), and on recall
both of those objectives are worse than MSE after correction, so their gain
in precision comes with a loss of coverage. Table~\ref{tab:full} lists the
mean and sample standard deviation of all five metrics for every objective of
the five-seed comparison, from the same runs as the paired tables.

% The tables follow in the order the text above cites them. Most are about
% two thirds of a page tall, too tall to share a page with text, so they are
% input after the prose rather than between its sentences, which would leave
% part-empty text pages. Each generated table keeps a placeholder that
% carries its label, as the section files do, so a missing file shows in the
% PDF rather than breaking every reference to it.
\IfFileExists{tables/price_curve.tex}{% Generated by scripts/make_tables.py --- do not edit by hand.
% Regenerate after every sweep; the paper reads these directly.

\begin{table}[tbp]
\centering
\small
\caption{The value of the high-band boost as a function of its dose, on MSE's samples (CIFAR-10, epoch 199, DDIM-100, five seeds, scored in floating point), and the prices on CelebA-64. FID is mean $\pm$ sample standard deviation; value $= M(B) - M(T(B))$ ($M$ is FID), paired by seed with its 95\% $t$-interval; the deficit is the high-band deficit after the operator (dB).}
\label{tab:price-curve}
\begin{tabular}{lrrr}
\toprule
Operator on MSE's samples & FID & Value [95\% CI] & Deficit after (dB) \\
\midrule
$+0.1$ dB & $18.446 \pm 0.211$ & $+0.014$\;$[0.012, 0.016]$ & $2.23 \pm 0.03$ \\
$+0.2$ dB & $18.432 \pm 0.211$ & $+0.029$\;$[0.026, 0.032]$ & $2.13 \pm 0.03$ \\
$+0.41$ dB & $18.399 \pm 0.211$ & $+0.061$\;$[0.056, 0.066]$ & $1.94 \pm 0.03$ \\
$+0.8$ dB & $18.335 \pm 0.212$ & $+0.126$\;$[0.117, 0.134]$ & $1.58 \pm 0.03$ \\
$+1.2$ dB & $18.270 \pm 0.211$ & $+0.190$\;$[0.182, 0.199]$ & $1.21 \pm 0.03$ \\
$+1.6$ dB & $18.214 \pm 0.206$ & $+0.246$\;$[0.237, 0.256]$ & $0.85 \pm 0.03$ \\
$+2.4$ dB & $18.147 \pm 0.201$ & $+0.313$\;$[0.297, 0.330]$ & $0.11 \pm 0.03$ \\
\midrule
\multicolumn{4}{l}{\emph{CelebA-64, MSE, epoch 199}} \\
nothing (as generated) & $6.195 \pm 0.049$ & --- & $2.33 \pm 0.02$ \\
boost, $+0.41$ dB & $6.223 \pm 0.050$ & $-0.028$\;$[-0.029, -0.027]$ & $1.95 \pm 0.02$ \\
boost, $+1.25$ dB & $6.324 \pm 0.053$ & $-0.130$\;$[-0.137, -0.123]$ & $1.18 \pm 0.02$ \\
full radial-spectrum match & $6.545 \pm 0.053$ & $-0.351$\;$[-0.362, -0.340]$ & $0.02 \pm 0.00$ \\
\bottomrule
\end{tabular}
\end{table}
}{%
  \begin{table}[tbp]\centering
  \fbox{\textbf{tables/price\_curve.tex not built}}
  \caption{Not built.}\label{tab:price-curve}
  \end{table}}
\IfFileExists{tables/operators.tex}{% Generated by scripts/make_tables.py --- do not edit by hand.
% Regenerate after every sweep; the paper reads these directly.

\begin{table}[tbp]
\centering
\small
\caption{Spectral operators applied to MSE's own samples (CIFAR-10, epoch 199, DDIM-100, five seeds). Unless marked, the filtered images are scored in floating point, without 8-bit quantisation. FID is mean $\pm$ sample standard deviation. The value of an operator is $M(B) - M(T(B))$, the FID it removes ($M$ is FID), paired by seed, with its 95\% $t$-interval. The deficit is the high-band deficit after the operator (dB). The transplant imposes, seed by seed, the radial profile of the AWWL run with the same seed; the match imposes the reference set's profile and so removes the mean radial deficit, not per-image or anisotropic differences.}
\label{tab:operators}
\begin{tabular}{lrrr}
\toprule
Applied to MSE's samples & FID & Value [95\% CI] & Deficit after (dB) \\
\midrule
nothing (as generated) & $18.460 \pm 0.210$ & --- & $2.32 \pm 0.03$ \\
high-band boost, $+0.41$~dB & $18.399 \pm 0.211$ & $+0.061$\;$[0.056, 0.066]$ & $1.94 \pm 0.03$ \\
\quad the same, rounded to 8 bits with dither & $18.119 \pm 0.201$ & $+0.342$\;$[0.323, 0.361]$ & $1.88 \pm 0.03$ \\
AWWL's radial profile transplanted & $18.118 \pm 0.145$ & $+0.343$\;$[0.241, 0.444]$ & $1.95 \pm 0.12$ \\
full radial-spectrum match & $17.759 \pm 0.193$ & $+0.701$\;$[0.672, 0.730]$ & $0.00 \pm 0.00$ \\
\quad the same, rounded to 8 bits & $17.537 \pm 0.192$ & $+0.923$\;$[0.886, 0.961]$ & $-0.04 \pm 0.01$ \\
\midrule
\emph{AWWL as trained} & $18.901 \pm 0.344$ & --- & $1.95 \pm 0.11$ \\
\bottomrule
\end{tabular}
\end{table}
}{%
  \begin{table}[tbp]\centering
  \fbox{\textbf{tables/operators.tex not built}}
  \caption{Not built.}\label{tab:operators}
  \end{table}}
\IfFileExists{tables/pricing_arms.tex}{% Generated by scripts/make_tables.py --- do not edit by hand.
% Regenerate after every sweep; the paper reads these directly.

\begin{table}[tbp]
\centering
\footnotesize
\setlength{\tabcolsep}{3pt}
\caption{The pricing test applied to every trained arm $A$ (CIFAR-10, epoch 199, DDIM-100). $T_A$ transplants the radial profile of $A$'s run onto the MSE run with the same seed, $B$; $\delta_{\mathrm{high}}$ is the resulting change in the high-band deviation, measured back on $T_A(B)$ (dB). With $M$ the FID, value $= M(B) - M(T_A(B))$, gap $= M(A) - M(B)$ and MRG $= M(A) - M(T_A(B))$, each a mean of paired-by-seed differences over the $n$ seeds with all three measurements; the MRG carries its 95\% $t$-interval (the gap's is in Table~\ref{tab:compare-fid}). $p$ is the two-sided paired $t$-test of MRG $= 0$, raw and Holm-corrected over the fixed family of the 19 arms listed: the 18 arms of the main family of Table~\ref{tab:compare-fid} that trained stably (Raw SNR weighting, Multi-level (3 levels) and GradNorm balancing did not and are not priced) and AWWL + EMA. $B$ is MSE without EMA for every arm, the two EMA arms included, so the MSE + EMA row measures weight averaging rather than an objective, and AWWL + EMA is priced against MSE, not against the MSE + EMA of its EMA-matched comparison in Table~\ref{tab:compare-fid}. Clipping to the pixel range makes $\delta_{\mathrm{high}}$ differ from the achieved change $s(A) - s(B)$ (Table~\ref{tab:bands} for the arms listed there; Appendix~\ref{app:scoring}) by at most $0.03$~dB, except for FFL ($+2.22$ against $+2.13$~dB). Bold: $p_{\mathrm{Holm}} < 0.05$ within the stated family.}
\label{tab:pricing-arms}
\begin{tabular}{lrrrrrrr}
\toprule
Arm $A$ & $n$ & $\delta_{\mathrm{high}}$ & Gap & Value & MRG [95\% CI] & $p$ & $p_{\mathrm{Holm}}$ \\
\midrule
\multicolumn{8}{l}{\emph{Conference-paper arms}} \\
AWWL & 5 & $+0.37$ & $+0.44$ & $+0.343$ & $+0.78$\;$[0.28, 1.29]$ & $0.013$ & $0.063$ \\
Static wavelet ($p = 0$) & 5 & $-0.04$ & $+0.17$ & $-0.018$ & $+0.16$\;$[-0.17, 0.48]$ & $0.25$ & $0.50$ \\
AWWL, detail bands summed & 5 & $+0.16$ & $+2.63$ & $+0.130$ & $+2.76$\;$[2.39, 3.13]$ & $<10^{-4}$ & $\mathbf{0.0005}$ \\
\addlinespace[2pt]
\multicolumn{8}{l}{\emph{Normalisation controls}} \\
AWWL, $\sum_b \omega_b = 1$ & 5 & $+0.29$ & $+1.15$ & $+0.271$ & $+1.42$\;$[0.80, 2.04]$ & $0.0031$ & $\mathbf{0.019}$ \\
AWWL, $\sum_b \omega_b = 0.5$ & 5 & $+0.26$ & $+1.34$ & $+0.250$ & $+1.59$\;$[1.29, 1.89]$ & $0.0001$ & $\mathbf{0.0014}$ \\
AWWL, $\sum_b \omega_b^2$ fixed & 5 & $+0.30$ & $+0.88$ & $+0.269$ & $+1.15$\;$[0.79, 1.52]$ & $0.0009$ & $\mathbf{0.0073}$ \\
AWWL, detail bands summed, $\sum_b \omega_b^2$ fixed & 5 & $+0.08$ & $+1.75$ & $+0.129$ & $+1.88$\;$[1.39, 2.38]$ & $0.0005$ & $\mathbf{0.0041}$ \\
\addlinespace[2pt]
\multicolumn{8}{l}{\emph{Other weightings}} \\
Min-SNR-$\gamma$ & 5 & $+0.32$ & $+2.09$ & $-0.056$ & $+2.03$\;$[1.74, 2.32]$ & $<10^{-4}$ & $\mathbf{0.0006}$ \\
FFL & 5 & $+2.22$ & $+2.24$ & $+0.420$ & $+2.66$\;$[2.21, 3.11]$ & $<10^{-4}$ & $\mathbf{0.0009}$ \\
\addlinespace[2pt]
\multicolumn{8}{l}{\emph{Pixel losses}} \\
$L_1$ & 5 & $-0.93$ & $+7.14$ & $-2.683$ & $+4.46$\;$[4.09, 4.83]$ & $<10^{-4}$ & $\mathbf{<10^{-4}}$ \\
Huber & 5 & $-0.06$ & $+0.28$ & $-0.043$ & $+0.24$\;$[-0.10, 0.58]$ & $0.12$ & $0.36$ \\
Charbonnier & 5 & $-0.95$ & $+7.08$ & $-2.689$ & $+4.39$\;$[4.14, 4.64]$ & $<10^{-4}$ & $\mathbf{<10^{-4}}$ \\
\addlinespace[2pt]
\multicolumn{8}{l}{\emph{Extensions (unnormalised)}} \\
Per-sub-band schedules & 5 & $+0.07$ & $+2.61$ & $+0.112$ & $+2.72$\;$[2.28, 3.17]$ & $<10^{-4}$ & $\mathbf{0.0009}$ \\
Spatially adaptive & 5 & $-0.44$ & $+4.65$ & $-0.199$ & $+4.45$\;$[3.98, 4.91]$ & $<10^{-4}$ & $\mathbf{0.0002}$ \\
Learned weighting & 5 & $+0.06$ & $+6.66$ & $-0.097$ & $+6.56$\;$[6.05, 7.08]$ & $<10^{-4}$ & $\mathbf{<10^{-4}}$ \\
Learned basis (lifting) & 3 & $+0.31$ & $+6.13$ & $+0.045$ & $+6.18$\;$[5.00, 7.36]$ & $0.0020$ & $\mathbf{0.014}$ \\
All axes combined & 3 & $-0.52$ & $+12.5$ & $-0.279$ & $+12.3$\;$[11.2, 13.3]$ & $0.0004$ & $\mathbf{0.0036}$ \\
\addlinespace[2pt]
\multicolumn{8}{l}{\emph{With EMA}} \\
MSE + EMA & 3 & $-0.04$ & $-0.03$ & $-0.011$ & $-0.04$\;$[-0.33, 0.24]$ & $0.58$ & $0.58$ \\
AWWL + EMA & 3 & $+0.38$ & $+0.45$ & $+0.338$ & $+0.79$\;$[-0.16, 1.74]$ & $0.069$ & $0.28$ \\
\addlinespace[2pt]
\multicolumn{8}{l}{\emph{Outside the declared family}} \\
AWWL, timestep only & 5 & $+0.14$ & $+0.81$ & $+0.045$ & $+0.85$\;$[0.53, 1.17]$ & $0.0018$ & --- \\
MSE, planted change & 5 & $+0.40$ & $-0.44$ & $+0.358$ & $-0.09$\;$[-0.25, 0.08]$ & $0.24$ & --- \\
\bottomrule
\end{tabular}
\end{table}
}{%
  \begin{table}[tbp]\centering
  \fbox{\textbf{tables/pricing\_arms.tex not built}}
  \caption{Not built.}\label{tab:pricing-arms}
  \end{table}}
\IfFileExists{tables/worth.tex}{% Generated by scripts/make_tables.py --- do not edit by hand.
% Regenerate after every sweep; the paper reads these directly.

\begin{table}[tbp]
\centering
\small
\caption{The high-band boost of $+0.41$~dB applied as a Fourier post-process to each arm's own samples (CIFAR-10, epoch 199, DDIM-100, \num{10000} samples per seed, five seeds), scored without 8-bit quantisation. The deficit is the high-band spectral deficit of the samples as generated (dB, positive means too little energy). FID is mean $\pm$ sample standard deviation. The value of the boost is the FID it removes, $M(\text{as generated}) - M(\text{boosted})$ ($M$ is FID), paired by seed, with its 95\% $t$-interval.}
\label{tab:worth}
\begin{tabular}{lrrrr}
\toprule
 & Deficit (dB) & FID & FID after boost & Value [95\% CI] \\
\midrule
MSE & $2.32 \pm 0.03$ & $18.460 \pm 0.210$ & $18.399 \pm 0.211$ & $+0.061$\;$[0.056, 0.066]$ \\
AWWL & $1.95 \pm 0.11$ & $18.901 \pm 0.344$ & $18.846 \pm 0.344$ & $+0.055$\;$[0.050, 0.061]$ \\
\bottomrule
\end{tabular}
\end{table}
}{%
  \begin{table}[tbp]\centering
  \fbox{\textbf{tables/worth.tex not built}}
  \caption{Not built.}\label{tab:worth}
  \end{table}}
\IfFileExists{tables/features.tex}{% Generated by scripts/make_tables.py --- do not edit by hand.
% Regenerate after every sweep; the paper reads these directly.

\begin{table}[tbp]
\centering
\small
\setlength{\tabcolsep}{4pt}
\caption{The pricing test in two feature spaces: FID (Inception-V3 pool features) and FD$_{\text{DINOv2}}$ (Fr\'echet distance in DINOv2 ViT-L/14 features), both scored in floating point on the same samples. CIFAR-10, epoch 199, DDIM-100, five seeds unless stated; value $= M(B) - M(T(B))$, gaps and MRG paired by seed with 95\% $t$-intervals (uncorrected); positive values are improvements for a value and costs for a gap or MRG. \textsuperscript{a}One measurement of the 10{,}000 CIFAR-10 test images against the training images; negative means the noise raises the distance. \textsuperscript{b}Three seeds of the 800-epoch MSE run.}
\label{tab:features}
\begin{tabular}{lrr}
\toprule
 & FID & FD$_{\text{DINOv2}}$ \\
\midrule
MSE, $M(B)$ (mean $\pm$ SD) & $18.460 \pm 0.210$ & $745.0 \pm 1.8$ \\
AWWL's gap to MSE, $D = M(A) - M(B)$ & $+0.44$\;$[-0.15, 1.03]$ & $+24.5$\;$[20.5, 28.5]$ \\
Value of AWWL's change, $V$ (transplant) & $+0.34$\;$[0.24, 0.44]$ & $-0.21$\;$[-0.48, 0.05]$ \\
$\mathrm{MRG} = M(A) - M(T_A(B))$ & $+0.78$\;$[0.28, 1.29]$ & $+24.3$\;$[20.2, 28.3]$ \\
\midrule
$V$, mid third only & $+0.38$\;$[0.25, 0.51]$ & $+1.03$\;$[0.55, 1.51]$ \\
$V$, high third only & $+0.03$\;$[0.02, 0.05]$ & $+0.18$\;$[0.004, 0.36]$ \\
Value of the $+0.41$~dB boost & $+0.06$\;$[0.06, 0.07]$ & $+0.41$\;$[0.30, 0.51]$ \\
Value of the full radial match & $+0.70$\;$[0.67, 0.73]$ & $+2.86$\;$[2.25, 3.46]$ \\
Value of white noise, $\pm0.5$~LSB & $+0.32$\;$[0.30, 0.33]$ & $-2.08$\;$[-2.49, -1.68]$ \\
Value of white noise, $\pm1$~LSB & $+0.93$\;$[0.87, 0.99]$ & $-6.92$\;$[-7.86, -5.97]$ \\
The same $\pm0.5$~LSB on held-out real images\textsuperscript{a} & $-0.01$ & $-0.14$ \\
\midrule
Positive control, MRG (epoch 399 vs.\ 199)\textsuperscript{b} & $-2.87$\;$[-3.50, -2.25]$ & $-118$\;$[-123, -114]$ \\
Planted-mechanism control, MRG & $-0.09$\;$[-0.25, 0.08]$ & $-0.69$\;$[-5.01, 3.63]$ \\
\midrule
AWWL, timestep only: gap to MSE & $+0.81$\;$[0.48, 1.14]$ & $+24.0$\;$[18.6, 29.5]$ \\
AWWL, $\sum_b w_b = 0.5$: gap to MSE & $+1.34$\;$[0.98, 1.70]$ & $+43.1$\;$[39.0, 47.2]$ \\
\midrule
FFL: value of its change, $V$ & $+0.42$\;$[0.39, 0.45]$ & $+0.70$\;$[0.23, 1.18]$ \\
FFL: MRG & $+2.66$\;$[2.21, 3.11]$ & $+87.1$\;$[81.8, 92.4]$ \\
\midrule
CelebA-64: value of the full radial match & $-0.35$\;$[-0.36, -0.34]$ & $-38.8$\;$[-40.2, -37.5]$ \\
CelebA-64: value of white noise, $\pm0.5$~LSB & $+0.42$\;$[0.42, 0.43]$ & $+6.95$\;$[6.65, 7.24]$ \\
\bottomrule
\end{tabular}
\end{table}
}{%
  \begin{table}[tbp]\centering
  \fbox{\textbf{tables/features.tex not built}}
  \caption{Not built.}\label{tab:features}
  \end{table}}
\IfFileExists{tables/compare_fid.tex}{% Generated by scripts/make_tables.py --- do not edit by hand.
% Regenerate after every sweep; the paper reads these directly.

\begin{table}[tbp]
\centering
\small
\setlength{\tabcolsep}{3pt}
\caption{Paired-by-seed comparison against MSE on FID, CIFAR-10, epoch 199, DDIM-100. $\Delta$ is the mean difference (objective $-$ MSE) over the $n$ shared training seeds, with its 95\% $t$-interval; $p$ is the two-sided paired $t$-test, raw and Holm--Bonferroni corrected over the fixed family of 21 arms (every row without a dagger). $^{\dagger}$A family of its own: AWWL + EMA is paired with MSE + EMA (one test); the timestep-only arm is corrected with AWWL against it, the two-test decomposition family (Section~\ref{sec:confound}). $p_{\mathrm{W}}$ is the two-sided Wilcoxon signed-rank test, whose smallest attainable value is $0.0625$ at five pairs and $0.25$ at three. Bold: $p_{\mathrm{Holm}} < 0.05$ within the stated family. Positive $\Delta$ is worse.}
\label{tab:compare-fid}
\begin{tabular}{lrrrrrr}
\toprule
Objective & $n$ & $\Delta$ & 95\% CI & $p$ & $p_{\mathrm{Holm}}$ & $p_{\mathrm{W}}$ \\
\midrule
\multicolumn{7}{l}{\emph{Conference-paper arms}} \\
AWWL & 5 & $+0.44$ & $[-0.15, 1.03]$ & $0.11$ & $0.32$ & $0.12$ \\
Static wavelet ($p = 0$) & 5 & $+0.17$ & $[-0.16, 0.51]$ & $0.22$ & $0.44$ & $0.44$ \\
AWWL, detail bands summed & 5 & $+2.63$ & $[2.22, 3.04]$ & $<10^{-4}$ & $\mathbf{0.0008}$ & $0.0625$ \\
\addlinespace[2pt]
\multicolumn{7}{l}{\emph{Normalisation controls}} \\
AWWL, $\sum_b \omega_b = 1$ & 5 & $+1.15$ & $[0.48, 1.82]$ & $0.0090$ & $\mathbf{0.045}$ & $0.0625$ \\
AWWL, $\sum_b \omega_b = 0.5$ & 5 & $+1.34$ & $[0.98, 1.70]$ & $0.0005$ & $\mathbf{0.0049}$ & $0.0625$ \\
AWWL, $\sum_b \omega_b^2$ fixed & 5 & $+0.88$ & $[0.47, 1.30]$ & $0.0041$ & $\mathbf{0.024}$ & $0.0625$ \\
AWWL, detail bands summed, $\sum_b \omega_b^2$ fixed & 5 & $+1.75$ & $[1.19, 2.32]$ & $0.0010$ & $\mathbf{0.0090}$ & $0.0625$ \\
\addlinespace[2pt]
\multicolumn{7}{l}{\emph{Timestep schedule alone}} \\
AWWL, timestep only & 5 & $+0.81$ & $[0.48, 1.14]$ & $0.0025$ & $\mathbf{0.0050}$$^{\dagger}$ & $0.0625$ \\
\addlinespace[2pt]
\multicolumn{7}{l}{\emph{Other weightings}} \\
Min-SNR-$\gamma$ & 5 & $+2.09$ & $[1.78, 2.39]$ & $<10^{-4}$ & $\mathbf{0.0007}$ & $0.0625$ \\
Raw SNR weighting & 5 & $+306$ & $[287, 325]$ & $<10^{-4}$ & $\mathbf{<10^{-4}}$ & $0.0625$ \\
FFL & 5 & $+2.24$ & $[1.80, 2.68]$ & $0.0001$ & $\mathbf{0.0018}$ & $0.0625$ \\
\addlinespace[2pt]
\multicolumn{7}{l}{\emph{Pixel losses}} \\
$L_1$ & 5 & $+7.14$ & $[6.79, 7.50]$ & $<10^{-4}$ & $\mathbf{<10^{-4}}$ & $0.0625$ \\
Huber & 5 & $+0.28$ & $[-0.05, 0.61]$ & $0.076$ & $0.31$ & $0.12$ \\
Charbonnier & 5 & $+7.08$ & $[6.76, 7.40]$ & $<10^{-4}$ & $\mathbf{<10^{-4}}$ & $0.0625$ \\
\addlinespace[2pt]
\multicolumn{7}{l}{\emph{Extensions (unnormalised)}} \\
Per-sub-band schedules & 5 & $+2.61$ & $[2.11, 3.11]$ & $0.0001$ & $\mathbf{0.0017}$ & $0.0625$ \\
Spatially adaptive & 5 & $+4.65$ & $[4.10, 5.19]$ & $<10^{-4}$ & $\mathbf{0.0003}$ & $0.0625$ \\
Multi-level (3 levels) & 5 & $+88.9$ & $[50.4, 127.3]$ & $0.0030$ & $\mathbf{0.021}$ & $0.0625$ \\
Learned weighting & 5 & $+6.66$ & $[6.12, 7.20]$ & $<10^{-4}$ & $\mathbf{<10^{-4}}$ & $0.0625$ \\
GradNorm balancing & 5 & $+330$ & $[296, 365]$ & $<10^{-4}$ & $\mathbf{0.0002}$ & $0.0625$ \\
Learned basis (lifting) & 3 & $+6.13$ & $[4.99, 7.28]$ & $0.0019$ & $\mathbf{0.015}$ & $0.25$ \\
All axes combined & 3 & $+12.5$ & $[11.5, 13.5]$ & $0.0004$ & $\mathbf{0.0039}$ & $0.25$ \\
\addlinespace[2pt]
\multicolumn{7}{l}{\emph{With EMA}} \\
MSE + EMA & 3 & $-0.03$ & $[-0.31, 0.24]$ & $0.65$ & $0.65$ & $0.75$ \\
AWWL + EMA vs.\ MSE + EMA & 3 & $+0.49$ & $[-0.91, 1.88]$ & $0.27$ & $0.27$$^{\dagger}$ & $0.25$ \\
\bottomrule
\end{tabular}
\end{table}
}{%
  \begin{table}[tbp]\centering
  \fbox{\textbf{tables/compare\_fid.tex not built}}
  \caption{Not built.}\label{tab:compare-fid}
  \end{table}}
\IfFileExists{tables/compare_other.tex}{% Generated by scripts/make_tables.py --- do not edit by hand.
% Regenerate after every sweep; the paper reads these directly.

\begin{table}[p]
\centering
\noindent\rotatebox{90}{\begin{minipage}{\textheight}%
\centering
\footnotesize
\setlength{\tabcolsep}{2.8pt}
\caption{Paired-by-seed comparison against MSE on Inception Score, KID, precision and recall, CIFAR-10, epoch 199, DDIM-100 (\num{10000} samples per run). $\Delta$ is the mean difference (objective $-$ MSE) over the shared training seeds, five unless a superscript gives the number; $p$ is the two-sided paired $t$-test, raw and Holm--Bonferroni corrected, within each metric, over the fixed family of 21 arms (every row whose $p_{\mathrm{Holm}}$ carries no dagger). $^{\dagger}$A family of its own: AWWL + EMA is paired with MSE + EMA (one test); the timestep-only arm is corrected with AWWL against it on the same metric, the two-test decomposition family of Table~\ref{tab:confound}. KID is from \texttt{torch-fidelity} (Appendix~\ref{app:supp-kid}). An upward arrow means higher is better, so a negative $\Delta$ is worse on IS, precision and recall and a positive one on KID. FID, with intervals, is in Table~\ref{tab:compare-fid}. Bold: $p_{\mathrm{Holm}} < 0.05$ within the stated family.}
\label{tab:compare-other}
\begin{tabular}{>{\raggedright\arraybackslash\hangindent=1em\hangafter=1}p{6.15cm}rrr@{\hspace{10pt}}rrr@{\hspace{10pt}}rrr@{\hspace{10pt}}rrr}
\toprule
 & \multicolumn{3}{c}{IS $\uparrow$} & \multicolumn{3}{c}{KID$\times 10^{3}$ $\downarrow$} & \multicolumn{3}{c}{Precision $\uparrow$} & \multicolumn{3}{c}{Recall $\uparrow$} \\
\cmidrule(lr){2-4}\cmidrule(lr){5-7}\cmidrule(lr){8-10}\cmidrule(lr){11-13}
Objective & $\Delta$ & $p$ & $p_{\mathrm{Holm}}$ & $\Delta$ & $p$ & $p_{\mathrm{Holm}}$ & $\Delta$ & $p$ & $p_{\mathrm{Holm}}$ & $\Delta$ & $p$ & $p_{\mathrm{Holm}}$ \\
\midrule
\multicolumn{13}{l}{\emph{Conference-paper arms}} \\
AWWL & $-0.059$ & $0.22$ & $1.00$ & $+0.31$ & $0.26$ & $1.00$ & $+0.0023$ & $0.29$ & $1.00$ & $-0.0068$ & $0.0094$ & $0.066$ \\
Static wavelet ($p = 0$) & $+0.013$ & $0.74$ & $1.00$ & $+0.30$ & $0.033$ & $0.17$ & $+0.0038$ & $0.22$ & $1.00$ & $-0.0005$ & $0.84$ & $1.00$ \\
AWWL, detail bands summed & $-0.235$ & $0.0019$ & $\mathbf{0.027}$ & $+2.72$ & $<10^{-4}$ & $\mathbf{0.0002}$ & $+0.0052$ & $0.12$ & $1.00$ & $-0.0181$ & $0.0025$ & $\mathbf{0.027}$ \\
\addlinespace[2pt]
\multicolumn{13}{l}{\emph{Normalisation controls}} \\
AWWL, $\sum_b \omega_b = 1$ & $-0.157$ & $0.0006$ & $\mathbf{0.0096}$ & $+1.31$ & $0.0031$ & $\mathbf{0.025}$ & $+0.0025$ & $0.32$ & $1.00$ & $-0.0095$ & $0.0066$ & $0.060$ \\
AWWL, $\sum_b \omega_b = 0.5$ & $-0.140$ & $0.0048$ & $0.062$ & $+1.47$ & $0.0005$ & $\mathbf{0.0057}$ & $-0.0005$ & $0.78$ & $1.00$ & $-0.0112$ & $0.0067$ & $0.060$ \\
AWWL, $\sum_b \omega_b^2$ fixed & $-0.118$ & $0.017$ & $0.15$ & $+0.98$ & $0.0007$ & $\mathbf{0.0064}$ & $+0.0012$ & $0.53$ & $1.00$ & $-0.0094$ & $0.037$ & $0.19$ \\
AWWL, detail bands summed, $\sum_b \omega_b^2$ fixed & $-0.198$ & $0.0017$ & $\mathbf{0.026}$ & $+2.31$ & $0.0002$ & $\mathbf{0.0026}$ & $+0.0034$ & $0.029$ & $0.44$ & $-0.0145$ & $0.010$ & $0.066$ \\
\addlinespace[2pt]
\multicolumn{13}{l}{\emph{Timestep schedule alone}} \\
AWWL, timestep only & $-0.006$ & $0.89$ & $0.89$$^{\dagger}$ & $+0.40$ & $0.028$ & $0.056$$^{\dagger}$ & $-0.0043$ & $0.026$ & $\mathbf{0.026}$$^{\dagger}$ & $+0.0046$ & $0.27$ & $0.27$$^{\dagger}$ \\
\addlinespace[2pt]
\multicolumn{13}{l}{\emph{Other weightings}} \\
Min-SNR-$\gamma$ & $-0.080$ & $0.060$ & $0.48$ & $-0.19$ & $0.41$ & $1.00$ & $+0.0067$ & $0.0010$ & $\mathbf{0.017}$ & $-0.0223$ & $0.0006$ & $\mathbf{0.0094}$ \\
Raw SNR weighting & $-6.428$ & $<10^{-4}$ & $\mathbf{<10^{-4}}$ & $+375$ & $<10^{-4}$ & $\mathbf{<10^{-4}}$ & $+0.0458$ & $0.67$ & $1.00$ & $-0.6410$ & $<10^{-4}$ & $\mathbf{<10^{-4}}$ \\
FFL & $-0.084$ & $0.068$ & $0.48$ & $-0.07$ & $0.53$ & $1.00$ & $-0.0145$ & $0.0002$ & $\mathbf{0.0049}$ & $+0.0020$ & $0.20$ & $0.60$ \\
\addlinespace[2pt]
\multicolumn{13}{l}{\emph{Pixel losses}} \\
$L_1$ & $-0.054$ & $0.30$ & $1.00$ & $+4.67$ & $<10^{-4}$ & $\mathbf{0.0001}$ & $-0.0161$ & $0.0007$ & $\mathbf{0.013}$ & $-0.0333$ & $<10^{-4}$ & $\mathbf{0.0004}$ \\
Huber & $+0.026$ & $0.44$ & $1.00$ & $+0.43$ & $0.010$ & $0.060$ & $+0.0002$ & $0.89$ & $1.00$ & $+0.0025$ & $0.075$ & $0.30$ \\
Charbonnier & $-0.035$ & $0.13$ & $0.81$ & $+4.49$ & $<10^{-4}$ & $\mathbf{<10^{-4}}$ & $-0.0140$ & $0.0031$ & $0.053$ & $-0.0329$ & $0.0003$ & $\mathbf{0.0058}$ \\
\addlinespace[2pt]
\multicolumn{13}{l}{\emph{Extensions (unnormalised)}} \\
Per-sub-band schedules & $-0.238$ & $0.0051$ & $0.062$ & $+2.53$ & $0.0001$ & $\mathbf{0.0017}$ & $+0.0027$ & $0.088$ & $1.00$ & $-0.0194$ & $0.0005$ & $\mathbf{0.0086}$ \\
Spatially adaptive & $-0.331$ & $0.0007$ & $\mathbf{0.011}$ & $+5.92$ & $<10^{-4}$ & $\mathbf{0.0001}$ & $-0.0020$ & $0.28$ & $1.00$ & $-0.0231$ & $0.0023$ & $\mathbf{0.027}$ \\
Multi-level (3 levels) & $-2.961$ & $0.0005$ & $\mathbf{0.0096}$ & $+64.8$ & $0.0069$ & $\mathbf{0.049}$ & $-0.0487$ & $0.053$ & $0.74$ & $-0.2811$ & $0.0013$ & $\mathbf{0.017}$ \\
Learned weighting & $-0.499$ & $0.0002$ & $\mathbf{0.0046}$ & $+5.84$ & $<10^{-4}$ & $\mathbf{0.0002}$ & $+0.0138$ & $<10^{-4}$ & $\mathbf{0.0012}$ & $-0.0576$ & $<10^{-4}$ & $\mathbf{0.0010}$ \\
GradNorm balancing & $-6.586$ & $<10^{-4}$ & $\mathbf{<10^{-4}}$ & $+378$ & $<10^{-4}$ & $\mathbf{0.0003}$ & $+0.0485$ & $0.65$ & $1.00$ & $-0.6410$ & $<10^{-4}$ & $\mathbf{<10^{-4}}$ \\
Learned basis (lifting)\,\textsuperscript{3} & $-0.537$ & $0.012$ & $0.12$ & $+6.71$ & $0.0005$ & $\mathbf{0.0057}$ & $-0.0015$ & $0.65$ & $1.00$ & $-0.0461$ & $0.0033$ & $\mathbf{0.033}$ \\
All axes combined\,\textsuperscript{3} & $-0.906$ & $0.0056$ & $0.062$ & $+13.6$ & $0.0003$ & $\mathbf{0.0030}$ & $+0.0189$ & $0.0061$ & $0.098$ & $-0.0993$ & $0.0007$ & $\mathbf{0.010}$ \\
\addlinespace[2pt]
\multicolumn{13}{l}{\emph{With EMA}} \\
MSE + EMA\,\textsuperscript{3} & $+0.001$ & $0.94$ & $1.00$ & $-0.02$ & $0.27$ & $1.00$ & $+0.0011$ & $0.35$ & $1.00$ & $-0.0005$ & $0.75$ & $1.00$ \\
AWWL + EMA\,\textsuperscript{3} vs.\ MSE + EMA & $-0.042$ & $0.55$ & $0.55$$^{\dagger}$ & $+0.28$ & $0.50$ & $0.50$$^{\dagger}$ & $-0.0012$ & $0.52$ & $0.52$$^{\dagger}$ & $-0.0058$ & $0.0082$ & $\mathbf{0.0082}$$^{\dagger}$ \\
\bottomrule
\end{tabular}
\end{minipage}}%
\end{table}
}{%
  \begin{table}[tbp]\centering
  \fbox{\textbf{tables/compare\_other.tex not built}}
  \caption{Not built.}\label{tab:compare-other}
  \end{table}}
\IfFileExists{tables/full.tex}{% Generated by scripts/make_tables.py --- do not edit by hand.
% Regenerate after every sweep; the paper reads these directly.

\begin{table}[tbp]
\centering
\footnotesize
\setlength{\tabcolsep}{3pt}
\caption{Every metric for every objective of the CIFAR-10 matrix: 200 epochs, epoch 199, DDIM-100, \num{10000} samples per run. Mean $\pm$ sample standard deviation over the $n$ training seeds. KID is from \texttt{torch-fidelity}; Appendix~\ref{app:supp-kid} explains why the hand-written estimator is not reported.}
\label{tab:full}
\begin{tabular}{>{\raggedright\arraybackslash\hangindent=1em\hangafter=1}p{4.6cm}rrrrrr}
\toprule
Objective & $n$ & FID $\downarrow$ & IS $\uparrow$ & KID$\times 10^{3}$ $\downarrow$ & Prec. $\uparrow$ & Rec. $\uparrow$ \\
\midrule
MSE & 5 & $18.46 \pm 0.21$ & $7.735 \pm 0.070$ & $13.34 \pm 0.27$ & $0.776 \pm 0.002$ & $0.641 \pm 0.003$ \\
\addlinespace[2pt]
\multicolumn{7}{l}{\emph{Conference-paper arms}} \\
AWWL & 5 & $18.90 \pm 0.34$ & $7.676 \pm 0.060$ & $13.65 \pm 0.45$ & $0.779 \pm 0.003$ & $0.634 \pm 0.003$ \\
Static wavelet ($p = 0$) & 5 & $18.63 \pm 0.06$ & $7.748 \pm 0.019$ & $13.64 \pm 0.16$ & $0.780 \pm 0.005$ & $0.640 \pm 0.004$ \\
AWWL, detail bands summed & 5 & $21.09 \pm 0.29$ & $7.500 \pm 0.023$ & $16.06 \pm 0.24$ & $0.782 \pm 0.006$ & $0.623 \pm 0.004$ \\
\addlinespace[2pt]
\multicolumn{7}{l}{\emph{Normalisation controls}} \\
AWWL, $\sum_b \omega_b = 1$ & 5 & $19.61 \pm 0.40$ & $7.578 \pm 0.035$ & $14.64 \pm 0.32$ & $0.779 \pm 0.005$ & $0.631 \pm 0.004$ \\
AWWL, $\sum_b \omega_b = 0.5$ & 5 & $19.80 \pm 0.19$ & $7.595 \pm 0.053$ & $14.80 \pm 0.17$ & $0.776 \pm 0.004$ & $0.630 \pm 0.004$ \\
AWWL, $\sum_b \omega_b^2$ fixed & 5 & $19.34 \pm 0.23$ & $7.617 \pm 0.098$ & $14.32 \pm 0.26$ & $0.778 \pm 0.003$ & $0.632 \pm 0.005$ \\
AWWL, detail bands summed, $\sum_b \omega_b^2$ fixed & 5 & $20.21 \pm 0.32$ & $7.536 \pm 0.058$ & $15.64 \pm 0.42$ & $0.780 \pm 0.003$ & $0.627 \pm 0.004$ \\
\addlinespace[2pt]
\multicolumn{7}{l}{\emph{Timestep schedule alone}} \\
AWWL, timestep only & 5 & $19.27 \pm 0.20$ & $7.729 \pm 0.095$ & $13.73 \pm 0.26$ & $0.772 \pm 0.002$ & $0.646 \pm 0.006$ \\
\addlinespace[2pt]
\multicolumn{7}{l}{\emph{Planted-mechanism control}} \\
MSE, planted change & 5 & $18.02 \pm 0.13$ & $7.824 \pm 0.030$ & $12.72 \pm 0.18$ & $0.772 \pm 0.002$ & $0.647 \pm 0.003$ \\
\addlinespace[2pt]
\multicolumn{7}{l}{\emph{Other weightings}} \\
Min-SNR-$\gamma$ & 5 & $20.55 \pm 0.10$ & $7.655 \pm 0.069$ & $13.14 \pm 0.24$ & $0.783 \pm 0.002$ & $0.619 \pm 0.003$ \\
Raw SNR weighting & 5 & $324.73 \pm 15.21$ & $1.307 \pm 0.017$ & $387.93 \pm 17.47$ & $0.822 \pm 0.222$ & $0.000 \pm 0.000$ \\
FFL & 5 & $20.70 \pm 0.18$ & $7.650 \pm 0.042$ & $13.27 \pm 0.25$ & $0.762 \pm 0.003$ & $0.643 \pm 0.003$ \\
\addlinespace[2pt]
\multicolumn{7}{l}{\emph{Pixel losses}} \\
$L_1$ & 5 & $25.60 \pm 0.17$ & $7.681 \pm 0.057$ & $18.01 \pm 0.20$ & $0.760 \pm 0.003$ & $0.608 \pm 0.004$ \\
Huber & 5 & $18.74 \pm 0.10$ & $7.761 \pm 0.039$ & $13.76 \pm 0.19$ & $0.777 \pm 0.002$ & $0.644 \pm 0.004$ \\
Charbonnier & 5 & $25.54 \pm 0.22$ & $7.700 \pm 0.077$ & $17.82 \pm 0.17$ & $0.762 \pm 0.003$ & $0.608 \pm 0.005$ \\
\addlinespace[2pt]
\multicolumn{7}{l}{\emph{Extensions (unnormalised)}} \\
Per-sub-band schedules & 5 & $21.07 \pm 0.27$ & $7.497 \pm 0.064$ & $15.87 \pm 0.20$ & $0.779 \pm 0.003$ & $0.622 \pm 0.004$ \\
Spatially adaptive & 5 & $23.11 \pm 0.35$ & $7.404 \pm 0.033$ & $19.26 \pm 0.33$ & $0.774 \pm 0.002$ & $0.618 \pm 0.005$ \\
Multi-level (3 levels) & 5 & $107.33 \pm 31.07$ & $4.774 \pm 0.689$ & $78.17 \pm 28.55$ & $0.728 \pm 0.041$ & $0.360 \pm 0.077$ \\
Learned weighting & 5 & $25.12 \pm 0.24$ & $7.236 \pm 0.033$ & $19.17 \pm 0.29$ & $0.790 \pm 0.003$ & $0.583 \pm 0.006$ \\
GradNorm balancing & 5 & $348.85 \pm 28.05$ & $1.149 \pm 0.066$ & $391.22 \pm 35.61$ & $0.825 \pm 0.217$ & $0.000 \pm 0.000$ \\
Learned basis (lifting) & 3 & $24.63 \pm 0.55$ & $7.187 \pm 0.029$ & $20.20 \pm 0.19$ & $0.775 \pm 0.005$ & $0.597 \pm 0.004$ \\
All axes combined & 3 & $31.03 \pm 0.12$ & $6.817 \pm 0.024$ & $27.08 \pm 0.31$ & $0.795 \pm 0.001$ & $0.544 \pm 0.004$ \\
\addlinespace[2pt]
\multicolumn{7}{l}{\emph{With EMA}} \\
MSE + EMA & 3 & $18.46 \pm 0.37$ & $7.724 \pm 0.079$ & $13.46 \pm 0.25$ & $0.778 \pm 0.001$ & $0.642 \pm 0.001$ \\
AWWL + EMA & 3 & $18.95 \pm 0.20$ & $7.683 \pm 0.065$ & $13.74 \pm 0.34$ & $0.776 \pm 0.002$ & $0.637 \pm 0.001$ \\
\bottomrule
\end{tabular}
\end{table}
}{%
  \begin{table}[tbp]\centering
  \fbox{\textbf{tables/full.tex not built}}
  \caption{Not built.}\label{tab:full}
  \end{table}}

\FloatBarrier
\section{Supplementary results}\label{app:supp}
This appendix gives the details, tables and figures that support the main
text but are not needed to follow it, in the order of the sections they
supplement. Appendices~\ref{app:supp-replication}--\ref{app:supp-extensions}
supplement the audit of Section~\ref{sec:audit} with the replication details,
the implementation dependence of KID, convergence, the sampler and training
budget, and the extensions of the weighting.
Appendix~\ref{app:supp-mechanism} repeats the band analysis of
Section~\ref{sec:mechanism} under DDPM-1000 and shows paired samples.
Appendix~\ref{app:supp-confound} reports a further normalisation control for
Section~\ref{sec:confound}. Appendices~\ref{app:supp-dreambooth}
and~\ref{sec:restoration} give the DreamBooth table and the
Gaussian-denoising experiment of Section~\ref{sec:generality}.

% Each file starts at subsection level and keeps the labels of the floats it
% carries; supp_generality.tex also carries tab:restoration, so no separate
% denoising subsection is input here.
\IfFileExists{sections/supp_audit.tex}{% Supplementary material for Section~\ref{sec:audit} (Auditing AWWL). Input
% under Appendix D (Supplementary results); starts at subsection level. Moved
% floats keep their labels (tab:replication-details is new): tab:convergence,
% fig:convergence, tab:sampler,
% tab:budget, fig:budget-sampler, fig:budget-curve, tab:deficit-budget,
% fig:deficit-metric, tab:extensions.

\subsection{Replication details}
\label{app:supp-replication}

Table~\ref{tab:replication-details} completes the FID comparison of
Table~\ref{tab:replication} with Inception Score, the one-sided tests against
the published differences and the 95\% prediction intervals for one further
paired run, under both samplers.

\IfFileExists{tables/replication_details.tex}{% Generated by scripts/make_tables.py --- do not edit by hand.
% Regenerate after every sweep; the paper reads these directly.

\begin{table}[tbp]
\centering
\footnotesize
\setlength{\tabcolsep}{3pt}
\caption{Replication details: every row of the CIFAR-10 table of the conference paper \citep{thao2026awwl} against the replication, on FID and Inception Score, epoch 199, \num{10000} samples per run. Published values are single runs (DDPM-1000, no EMA). Replicated values come from five training seeds per arm, the same weights sampled with DDPM-1000, as published, and with DDIM-100; only MSE and the two wavelet arms were sampled with DDPM-1000. Raw SNR weighting is the SNR-weighted row: the released code's unclamped weighting, which did not train stably in any seed. The Perceptual (VGG) row was not retrained (Appendix~\ref{app:setup-arms}). Differences are paired by seed, with 95\% $t$-intervals and two-sided paired $t$-tests, Holm-corrected over the two wavelet arms for DDPM-1000 (as in Table~\ref{tab:replication}) and over the 21-arm family of Table~\ref{tab:compare-fid} for DDIM-100. Bold: $p_{\mathrm{Holm}} < 0.05$ within the stated family. The one-sided test asks whether the published difference is the arm's mean effect: its $p$ is that of the alternative that the true difference is less favourable than published (larger for FID, smaller for Inception Score), uncorrected. The prediction interval asks whether it could be one run: it is the 95\% interval for the paired difference that one further run would show. A one-sided $p$ near 1 says only that the replicated difference is not less favourable than published; a replication more favourable than published (Huber on FID) shows in the prediction interval instead. In the lower three panels the Published columns repeat the published difference.}
\label{tab:replication-details}
\begin{tabular}{>{\raggedright\arraybackslash\hangindent=1em\hangafter=1}p{3.2cm}rrr@{\hspace{9pt}}rrr}
\toprule
 & \multicolumn{3}{c}{FID $\downarrow$} & \multicolumn{3}{c}{Inception Score $\uparrow$} \\
\cmidrule(lr){2-4}\cmidrule(lr){5-7}
 & Published & DDPM-1000 & DDIM-100 & Published & DDPM-1000 & DDIM-100 \\
\midrule
\multicolumn{7}{l}{\emph{Published single run; replicated mean $\pm$ sample standard deviation over five seeds}} \\
MSE & $16.68$ & $16.54 \pm 0.16$ & $18.46 \pm 0.21$ & $7.78$ & $7.834 \pm 0.055$ & $7.735 \pm 0.070$ \\
$L_1$ & $19.70$ & not run & $25.60 \pm 0.17$ & $7.69$ & not run & $7.681 \pm 0.057$ \\
Huber & $20.13$ & not run & $18.74 \pm 0.10$ & $7.62$ & not run & $7.761 \pm 0.039$ \\
Perceptual (VGG) & $20.02$ & not run & not run & $7.70$ & not run & not run \\
Raw SNR weighting & $19.83$ & not run & $325 \pm 15$ & $7.58$ & not run & $1.307 \pm 0.017$ \\
Static wavelet ($p = 0$) & $16.55$ & $16.85 \pm 0.09$ & $18.63 \pm 0.06$ & $7.81$ & $7.856 \pm 0.056$ & $7.748 \pm 0.019$ \\
AWWL & $16.62$ & $17.52 \pm 0.14$ & $18.90 \pm 0.34$ & $7.95$ & $7.790 \pm 0.069$ & $7.676 \pm 0.060$ \\
\midrule
\multicolumn{7}{l}{\emph{Difference from MSE, paired by seed, with its 95\% $t$-interval}} \\
$L_1$ & $+3.02$ & --- & $+7.14$ & $-0.09$ & --- & $-0.054$ \\
\quad 95\% CI &  & --- & $[6.79, 7.50]$ &  & --- & $[-0.181, 0.073]$ \\
\quad $p$ (raw; Holm) &  & --- & $<10^{-4}$; $\mathbf{<10^{-4}}$ &  & --- & $0.30$; $1.00$ \\
Huber & $+3.45$ & --- & $+0.28$ & $-0.16$ & --- & $+0.026$ \\
\quad 95\% CI &  & --- & $[-0.05, 0.61]$ &  & --- & $[-0.058, 0.109]$ \\
\quad $p$ (raw; Holm) &  & --- & $0.076$; $0.31$ &  & --- & $0.44$; $1.00$ \\
Raw SNR weighting & $+3.15$ & --- & $+306$ & $-0.20$ & --- & $-6.428$ \\
\quad 95\% CI &  & --- & $[287, 325]$ &  & --- & $[-6.525, -6.331]$ \\
\quad $p$ (raw; Holm) &  & --- & $<10^{-4}$; $\mathbf{<10^{-4}}$ &  & --- & $<10^{-4}$; $\mathbf{<10^{-4}}$ \\
Static wavelet ($p = 0$) & $-0.13$ & $+0.31$ & $+0.17$ & $+0.03$ & $+0.022$ & $+0.013$ \\
\quad 95\% CI &  & $[0.19, 0.43]$ & $[-0.16, 0.51]$ &  & $[-0.038, 0.083]$ & $[-0.085, 0.110]$ \\
\quad $p$ (raw; Holm) &  & $0.0022$; $\mathbf{0.0022}$ & $0.22$; $0.44$ &  & $0.36$; $0.73$ & $0.74$; $1.00$ \\
AWWL & $-0.06$ & $+0.98$ & $+0.44$ & $+0.17$ & $-0.043$ & $-0.059$ \\
\quad 95\% CI &  & $[0.80, 1.16]$ & $[-0.15, 1.03]$ &  & $[-0.191, 0.104]$ & $[-0.171, 0.054]$ \\
\quad $p$ (raw; Holm) &  & $0.0001$; $\mathbf{0.0002}$ & $0.11$; $0.32$ &  & $0.46$; $0.73$ & $0.22$; $1.00$ \\
\midrule
\multicolumn{7}{l}{\emph{One-sided test of the published difference from MSE, $p$}} \\
$L_1$ & $+3.02$ & --- & $<10^{-4}$ & $-0.09$ & --- & $0.76$ \\
Huber & $+3.45$ & --- & $1.00$ & $-0.16$ & --- & $1.00$ \\
Raw SNR weighting & $+3.15$ & --- & $<10^{-4}$ & $-0.20$ & --- & $<10^{-4}$ \\
Static wavelet ($p = 0$) & $-0.13$ & $0.0003$ & $0.032$ & $+0.03$ & $0.37$ & $0.32$ \\
AWWL & $-0.06$ & $<10^{-4}$ & $0.039$ & $+0.17$ & $0.0080$ & $0.0024$ \\
\midrule
\multicolumn{7}{l}{\emph{95\% prediction interval for the difference from MSE of one further paired run}} \\
$L_1$ & $+3.02$ & --- & $[6.28, 8.01]$ & $-0.09$ & --- & $[-0.366, 0.258]$ \\
Huber & $+3.45$ & --- & $[-0.53, 1.09]$ & $-0.16$ & --- & $[-0.179, 0.230]$ \\
Raw SNR weighting & $+3.15$ & --- & $[260, 352]$ & $-0.20$ & --- & $[-6.665, -6.190]$ \\
Static wavelet ($p = 0$) & $-0.13$ & $[0.01, 0.60]$ & $[-0.64, 0.99]$ & $+0.03$ & $[-0.126, 0.171]$ & $[-0.227, 0.252]$ \\
AWWL & $-0.06$ & $[0.54, 1.42]$ & $[-1.01, 1.89]$ & $+0.17$ & $[-0.405, 0.319]$ & $[-0.334, 0.217]$ \\
\bottomrule
\end{tabular}
\end{table}
}{%
  \begin{table}[tbp]\centering
  \fbox{\textbf{tables/replication\_details.tex not built}}
  \caption{Not built.}\label{tab:replication-details}
  \end{table}}

On Inception Score neither wavelet objective differs from MSE under either
sampler, even before correction: AWWL's paired difference is $\ddpmDISAWWL$
$[\ddpmDISAWWLLo, \ddpmDISAWWLHi]$ under DDPM-1000 (raw $p = \ddpmDISAWWLP$)
and $\dISAWWL$ $[\dISAWWLLo, \dISAWWLHi]$ under DDIM-100 (raw
$p = \dISAWWLP$), and the static wavelet arm's differences lie nearer zero.

A published single-run difference can be read as the objective's mean effect
or as the outcome of one run. A one-sided test against the published
difference, uncorrected and under the published sampler, rejects the first
reading for AWWL on both metrics (FID: $p \pPubFIDAWWLDdpm$; Inception Score:
$p = \pPubISAWWLDdpm$) and for the static wavelet arm on FID only
($p = \pPubFIDStaticDdpm$, against $p = \pPubISStaticDdpm$ for Inception
Score). The published differences were measured with DDPM-1000; compared with
DDIM-100 means, which assumes that the published effect does not depend on
the sampler, the same tests give $p = \pPubFIDAWWLDdim$ (FID) and
$p = \pPubISAWWLDdim$ (Inception Score) for AWWL.

The second reading is judged against the prediction interval. AWWL's
published Inception Score gain of $+\pubDISAWWL$ lies inside it under both
samplers, so one run can show that gain through seed variation alone. The
published FID differences lie outside it under DDPM-1000 ($\pubDFIDAWWL$
against $[\piFIDAWWLDdpmLo, \piFIDAWWLDdpmHi]$ for AWWL, $\pubDFIDStatic$
against $[\piFIDStaticDdpmLo, \piFIDStaticDdpmHi]$ for the static wavelet
arm) and inside the wider intervals of DDIM-100. Under the published sampler,
seed variation can therefore account for the published Inception Score gain
but not for the published FID differences.

\subsection{Implementation dependence of KID}
\label{app:supp-kid}

The sign of AWWL's KID difference from MSE depends on the implementation.
The library used for KID in the comparison tables
(Tables~\ref{tab:compare-other} and~\ref{tab:full}), \texttt{torch-fidelity},
gives a difference of
$+\dKIDAWWL \times 10^{-3}$ ($[\dKIDAWWLLo, \dKIDAWWLHi] \times 10^{-3}$,
uncorrected $p = \dKIDAWWLP$), whereas a
hand-written implementation with the same kernel and unbiased estimator gives
a nominally significant difference in AWWL's favour
($\dKIDhandAWWL \times 10^{-3}$,
$[\dKIDhandAWWLLo, \dKIDhandAWWLHi] \times 10^{-3}$, uncorrected
$p = \dKIDhandAWWLP$). The two differ only in choices that the definition of
KID leaves open: the Inception weights and preprocessing, a reference subset
of \num{10000} training images rather than all \num{\nRefImages}, and one
estimate over all samples rather than an average over subsets. The
\texttt{clean-fid} implementation used for the pricing test
(Section~\ref{sec:worth-headline}) gives $+\kidGapAWWL \times 10^{-3}$
($[\kidGapAWWLLo, \kidGapAWWLHi] \times 10^{-3}$, uncorrected
$p = \kidGapAWWLP$). An evaluation with the hand-written implementation alone
would have reported a gain that neither library reproduces.

\FloatBarrier
\subsection{Convergence within the 200-epoch schedule}
\label{app:supp-convergence}

The intermediate checkpoints give no evidence that AWWL reaches a given
quality with less training (Table~\ref{tab:convergence} and
Fig.~\ref{fig:convergence}). Under DDIM-100 it trails MSE at epochs 39, 79,
119 and 199 and leads only at epoch 159 (AWWL $-$ MSE $= \convDOneFiftyNine$~FID
$[\convDOneFiftyNineLo, \convDOneFiftyNineHi]$, uncorrected
$p = \convDOneFiftyNineP$, lower for \convAheadOneFiftyNine{} of the five
seeds), where seed variation is much larger than at the final checkpoint (MSE
standard deviation $\convSDMSEOneFiftyNine$~FID against
$\convSDMSEOneNinetyNine$). No checkpoint shows a significant difference in
AWWL's favour. Under the cosine learning-rate schedule these checkpoints
trace one run rather than shorter training budgets;
Section~\ref{sec:budget} and Appendix~\ref{app:supp-budget} report a longer
budget.

\IfFileExists{tables/convergence.tex}{% Generated by scripts/make_tables.py --- do not edit by hand.
% Regenerate after every sweep; the paper reads these directly.

\begin{table}[tbp]
\centering
\small
\caption{FID through training, CIFAR-10, DDIM-100, five seeds. Mean $\pm$ sample standard deviation over seeds; $\Delta$ is AWWL $-$ MSE paired by seed with its 95\% $t$-interval and two-sided paired $t$-test $p$, uncorrected. The last column counts the seeds at which AWWL has the lower FID.}
\label{tab:convergence}
\begin{tabular}{rrrrrr}
\toprule
Epoch & MSE & AWWL & $\Delta$ [95\% CI] & $p$ & AWWL ahead \\
\midrule
39 & $52.38 \pm 6.07$ & $59.91 \pm 6.13$ & $+7.53$\;$[2.97, 12.09]$ & $0.010$ & 0 of 5 \\
79 & $33.91 \pm 1.53$ & $36.91 \pm 2.61$ & $+3.01$\;$[-0.15, 6.16]$ & $0.057$ & 1 of 5 \\
119 & $24.93 \pm 1.19$ & $26.43 \pm 0.82$ & $+1.50$\;$[0.12, 2.88]$ & $0.039$ & 0 of 5 \\
159 & $22.45 \pm 3.14$ & $21.18 \pm 1.32$ & $-1.27$\;$[-4.15, 1.62]$ & $0.29$ & 4 of 5 \\
199 & $18.46 \pm 0.21$ & $18.90 \pm 0.34$ & $+0.44$\;$[-0.15, 1.03]$ & $0.11$ & 1 of 5 \\
\bottomrule
\end{tabular}
\end{table}
}{%
  \begin{table}[tbp]\centering
  \fbox{\textbf{tables/convergence.tex not built}}
  \caption{Not built.}\label{tab:convergence}
  \end{table}}

\guardedfigure{1.0}{convergence.png}{%
  FID through training for MSE, the static wavelet arm and AWWL: CIFAR-10,
  DDIM-100, \num{10000} samples, no EMA, five seeds, evaluated at epochs 39,
  79, 119, 159 and 199. (a)~Means over seeds. (b)~Mean paired
  difference from MSE with its 95\% $t$-interval, the error bars offset
  horizontally for legibility; markers are hollow as no correction is applied
  across checkpoints. Table~\ref{tab:convergence} lists the values
  for AWWL.}{fig:convergence}

\FloatBarrier
\subsection{Sampler and training budget}
\label{app:supp-budget}

Table~\ref{tab:sampler} gives the DDPM-1000 resampling of
Section~\ref{sec:sampler} with its KID differences; its FID entries repeat
those of Table~\ref{tab:replication}.

\IfFileExists{tables/sampler.tex}{% Generated by scripts/make_tables.py --- do not edit by hand.
% Regenerate after every sweep; the paper reads these directly.

\begin{table}[tbp]
\centering
\small
\setlength{\tabcolsep}{3pt}
\caption{The same trained weights under two samplers, five seeds, \num{10000} samples per run, fixed sampling seed. FID is mean $\pm$ sample standard deviation. $\Delta$ (objective $-$ MSE, paired by seed, with its 95\% $t$-interval) and $p$ refer to DDPM-1000; Holm correction is over the two arms shown, separately for FID and KID (KID from \texttt{torch-fidelity}, $\times 10^{3}$). Bold: $p_{\mathrm{Holm}} < 0.05$ within the stated family.}
\label{tab:sampler}
\begin{tabular}{lrrrrrrr}
\toprule
 & \multicolumn{2}{c}{FID} & \multicolumn{3}{c}{FID, DDPM-1000} & \multicolumn{2}{c}{KID, DDPM-1000} \\
\cmidrule(lr){2-3}\cmidrule(lr){4-6}\cmidrule(lr){7-8}
Objective & DDIM-100 & DDPM-1000 & $\Delta$ [95\% CI] & $p$ & $p_{\mathrm{Holm}}$ & $\Delta$ & $p_{\mathrm{Holm}}$ \\
\midrule
MSE & $18.46 \pm 0.21$ & $16.54 \pm 0.16$ & --- & --- & --- & --- & --- \\
Static wavelet ($p = 0$) & $18.63 \pm 0.06$ & $16.85 \pm 0.09$ & $+0.31$\;$[0.19, 0.43]$ & $0.0022$ & $\mathbf{0.0022}$ & $+0.31$ & $0.054$ \\
AWWL & $18.90 \pm 0.34$ & $17.52 \pm 0.14$ & $+0.98$\;$[0.80, 1.16]$ & $0.0001$ & $\mathbf{0.0002}$ & $+1.09$ & $\mathbf{0.0005}$ \\
\bottomrule
\end{tabular}
\end{table}
}{%
  \begin{table}[tbp]\centering
  \fbox{\textbf{tables/sampler.tex not built}}
  \caption{Not built.}\label{tab:sampler}
  \end{table}}

Table~\ref{tab:budget} and Fig.~\ref{fig:budget-sampler} set AWWL's gap to MSE
at the two budgets of Section~\ref{sec:budget} side by side. The 800-epoch
runs differ from the main matrix in three respects: the budget, an
exponential moving average of the weights, and a cosine schedule stretched to
the new length. Within each budget the two objectives are trained
identically, so the table is read row by row. At 200 epochs EMA shows no
effect that could account for the difference between the rows, although three
seeds cannot exclude one: it leaves the FID of MSE where it was ($\dMSEema$
$[\dMSEemaLo, \dMSEemaHi]$) and AWWL's gap close to its value without EMA
($+\dEMApair$, $p = \dEMApairP$, against $+\dAWWL$). Of the gaps
in Table~\ref{tab:budget}, only the 200-epoch DDPM-1000 gap is significant.

\IfFileExists{tables/budget.tex}{% Generated by scripts/make_tables.py --- do not edit by hand.
% Regenerate after every sweep; the paper reads these directly.

\begin{table}[tbp]
\centering
\small
\setlength{\tabcolsep}{3pt}
\caption{AWWL against MSE at two training budgets and two samplers; $n$ is the number of seeds. FID is mean $\pm$ sample standard deviation over seeds. Each row is one operating point, and the comparison within a row is the meaningful one: $\Delta$ is the FID gap AWWL $-$ MSE paired by seed with its 95\% $t$-interval; $p$ is the two-sided paired $t$-test, raw and Holm-corrected, for the 800-epoch rows over the two 800-epoch operating points. $^{\dagger}$The 200-epoch rows repeat tests of other families and carry those families' corrected $p$: the main family of Table~\ref{tab:compare-fid} (no EMA, DDIM-100), the one-test EMA family (EMA, DDIM-100) and the sampler family of Table~\ref{tab:replication} (DDPM-1000). Bold: $p_{\mathrm{Holm}} < 0.05$ within the stated family. The 800-epoch runs use EMA and a cosine schedule stretched to their length.}
\label{tab:budget}
\begin{tabular}{llrrrrrr}
\toprule
Budget & Sampler & $n$ & MSE & AWWL & $\Delta$ [95\% CI] & $p$ & $p_{\mathrm{Holm}}$ \\
\midrule
200 epochs, no EMA & DDIM-100 & 5 & $18.46 \pm 0.21$ & $18.90 \pm 0.34$ & $+0.44$\;$[-0.15, 1.03]$ & $0.11$ & $0.32$$^{\dagger}$ \\
200 epochs, EMA & DDIM-100 & 3 & $18.46 \pm 0.37$ & $18.95 \pm 0.20$ & $+0.49$\;$[-0.91, 1.88]$ & $0.27$ & $0.27$$^{\dagger}$ \\
800 epochs, EMA & DDIM-100 & 3 & $13.70 \pm 0.15$ & $13.85 \pm 0.10$ & $+0.15$\;$[-0.44, 0.73]$ & $0.40$ & $0.40$ \\
200 epochs, no EMA & DDPM-1000 & 5 & $16.54 \pm 0.16$ & $17.52 \pm 0.14$ & $+0.98$\;$[0.80, 1.16]$ & $0.0001$ & $\mathbf{0.0002}$$^{\dagger}$ \\
800 epochs, EMA & DDPM-1000 & 3 & $11.96 \pm 0.10$ & $12.30 \pm 0.16$ & $+0.34$\;$[-0.22, 0.91]$ & $0.12$ & $0.24$ \\
\bottomrule
\end{tabular}
\end{table}
}{%
  \begin{table}[tbp]\centering
  \fbox{\textbf{tables/budget.tex not built}}
  \caption{Not built.}\label{tab:budget}
  \end{table}}

\guardedfigure{0.72}{budget_sampler.png}{%
  AWWL's paired FID gap to MSE at the two budgets of Table~\ref{tab:budget},
  under DDIM-100 (solid line, circles) and DDPM-1000 (dashed line, squares),
  with 95\% $t$-intervals. Filled markers: $p < 0.05$ after Holm correction
  within the contrast's family (Table~\ref{tab:families}). The 800-epoch
  points rest on three seeds and differ from the 200-epoch points in EMA and
  schedule length as well as in budget.}{fig:budget-sampler}

Fig.~\ref{fig:budget-curve} follows the longer runs through training. The FID
of MSE is lowest at epoch~\longBestEpochMSE{} (\longBestFIDMSE) and rises
afterwards, so the 800-epoch endpoint lies past the optimum of this recipe;
the recipe uses no dropout, and the rise is consistent with overfitting,
which we did not test. The paired gap is positive at every evaluated
checkpoint, between $+\longGapSixNinetyNine$ and $+\longGapOneNinetyNine$,
and no checkpoint's interval excludes zero. Because the checkpoints share the
same three seed pairs, they are repeated measurements of one set of runs, not
independent replications. Averaged over the checkpoints within each seed, the
gap is $+\longGapAvg$ $[\longGapAvgLo, \longGapAvgHi]$ ($p = \longGapAvgP$).

\guardedfigure{1.0}{budget_curve.png}{%
  The 800-epoch runs at every evaluated checkpoint (epochs 199 to 799 in steps
  of 100): EMA, stretched cosine schedule, DDIM-100, three seeds.
  (a)~Seed means of FID for MSE and AWWL. (b)~Mean paired
  difference AWWL $-$ MSE, with its 95\% $t$-interval shaded. The seven
  checkpoints share the same three seed pairs.}{fig:budget-curve}

The under-training reading predicts that the baseline's high-band deficit
(Section~\ref{sec:pricing-test}) shrinks as training continues;
Table~\ref{tab:deficit-budget} tests this. Compared end against end, each run
at the annealed end of its own schedule and both with EMA, the deficit of MSE
is larger after 800 epochs than after 200, by $+\trajEndToEnd$~dB
$[\trajEndToEndLo, \trajEndToEndHi]$, while its FID improves from
\budFIDMSEShortEmaDdim{} to \budFIDMSELongDdim. Within the 800-epoch run the
deficit rises from \trajDefMSEOneNinetyNine~dB at epoch~199 to
\trajDefMSESevenNinetyNine~dB at epoch~799 (Table~\ref{tab:deficit-budget} and
Fig.~\ref{fig:deficit-metric}), but these checkpoints differ in learning-rate
phase as well as in training length, the early ones being taken at a high
learning rate, so this trajectory is descriptive only. At
epoch~\longBestEpochMSE, where FID is lowest, the deficit does not differ from
that at the end of the 200-epoch schedule (Table~\ref{tab:deficit-budget},
last column). Over this range, more training improves FID without closing the
deficit. AWWL's high-band change at the end of the longer schedule
($+\trajCorrSevenNinetyNine$~dB
$[\trajCorrSevenNinetyNineLo, \trajCorrSevenNinetyNineHi]$) has a point
estimate below its value at the end of the 200-epoch schedule with EMA
($+\trajCorrEnd$~dB $[\trajCorrEndLo, \trajCorrEndHi]$), but the two
intervals overlap; Section~\ref{sec:worth} prices the change at this budget.

\IfFileExists{tables/deficit_budget.tex}{% Generated by scripts/make_band_table.py --- do not edit by hand.
% Recomputed from the sampled images (all samples, all reference images).

\begin{table}[tbp]
\centering
\footnotesize
\setlength{\tabcolsep}{2.5pt}
\caption{High-band spectral deficit (dB; minus the signed high-band deviation of Table~\ref{tab:bands}, so larger means less high-frequency energy) and FID against training budget, all with EMA, three seeds, DDIM-100 samples. The first row is the annealed end of the 200-epoch schedule; the others are checkpoints of the 800-epoch schedule, whose early checkpoints are taken at a high learning rate. Mean $\pm$ sample standard deviation over seeds. AWWL's correction is MSE's deficit minus AWWL's, paired by seed. The last two columns give the change in MSE's deficit, paired by seed, against epoch 199 of the same run and against the end of the 200-epoch schedule (first row). Brackets are 95\% $t$-intervals.}
\label{tab:deficit-budget}
\begin{tabular}{lrrrrrr}
\toprule
 & MSE & \multicolumn{2}{c}{Deficit (dB)} & AWWL & \multicolumn{2}{c}{MSE change in deficit} \\
\cmidrule(lr){3-4}\cmidrule(lr){6-7}
Checkpoint & FID & MSE & AWWL & correction & vs.\ epoch 199 & vs.\ 200 ep. \\
\midrule
200 ep., end (199) & $18.46 \pm 0.37$ & $2.36 \pm 0.02$ & $1.95 \pm 0.09$ & $+0.41$\;$[0.14, 0.67]$ & --- & --- \\
800 ep., epoch 199 & $14.78 \pm 0.19$ & $1.85 \pm 0.03$ & $1.67 \pm 0.04$ & $+0.18$\;$[0.04, 0.33]$ & --- & $-0.51$\;$[-0.59, -0.44]$ \\
800 ep., epoch 399 & $12.01 \pm 0.17$ & $2.31 \pm 0.02$ & $1.89 \pm 0.02$ & $+0.42$\;$[0.33, 0.51]$ & $+0.46$\;$[0.31, 0.60]$ & $-0.06$\;$[-0.13, 0.02]$ \\
800 ep., end (799) & $13.70 \pm 0.15$ & $2.60 \pm 0.02$ & $2.36 \pm 0.02$ & $+0.24$\;$[0.14, 0.34]$ & $+0.75$\;$[0.68, 0.81]$ & $+0.23$\;$[0.22, 0.25]$ \\
\bottomrule
\end{tabular}
\end{table}
}{%
  \begin{table}[tbp]\centering
  \fbox{\textbf{tables/deficit\_budget.tex not built}}
  \caption{Not built.}\label{tab:deficit-budget}
  \end{table}}

\guardedfigure{0.82}{deficit_vs_metric.png}{%
  High-band deficit and FID of MSE and AWWL through the 800-epoch runs.
  (a)~High-band spectral deficit of the samples (dB below the real
  images) and (b)~FID at epochs 199, 399 and 799 (EMA, DDIM-100), mean $\pm$ one sample standard
  deviation over three seeds. The checkpoints belong to one stretched cosine
  schedule and differ in learning-rate phase as well as in training length,
  so the figure is descriptive; Table~\ref{tab:deficit-budget} compares the
  annealed endpoints of the two schedules.}{fig:deficit-metric}

\FloatBarrier
\subsection{Extensions of the weighting}
\label{app:supp-extensions}

Table~\ref{tab:extensions} gives every extension of
Section~\ref{sec:extensions} against MSE and against its parent. Apart from
GradNorm, whose global weights sum to four, no arm fixes its total weight
over $\sigma_t$. The per-sub-band, spatially adaptive and learned-basis arms
give each detail band the full $w_{\det}$, so their parent is AWWL with detail
bands summed (Eq.~\eqref{eq:summed}); the three-level arm divides $w_{\det}$
among its nine detail bands, and the remaining arms set their own weights.

The extensions do not deliver the mechanism either
(Table~\ref{tab:pricing-arms}). The per-sub-band and learned-weighting arms
leave the high band close to that of MSE, the learned basis raises it by about
as much as AWWL, and the spatially adaptive and all-axes arms lower it.
Priced by transplant (Section~\ref{sec:pricing-test}), the last two changes
are worth $\valueArmSpatial$ and $\valueArmAllAxes$~FID: each raises the FID
of the MSE samples.

The learned weighting is uncertainty weighting \citep{kendall2018multi}
conditioned on the noise level: a small network maps $\sigma_t$ to a
log-variance $\psi_b$ for each band $b$, the band is weighted by
$\tfrac{1}{2}e^{-\psi_b}$, and the total is left free (a per-band form of the
EDM2 weighting \citep{karras2024analyzing}). Each weight is then driven
towards a value inversely proportional to the loss of its band at that noise
level, which produces the steep timestep schedule reported in
Section~\ref{sec:extensions}.

Two arms failed rather than degraded, and neither is priced. GradNorm
\citep{chen2018gradnorm}, run at asymmetry $1.5$ with global weights that sum
to four, drove the LL, LH and HL weights to their positivity floor in all five
seeds, training the network through the HH band alone; it ends at
\num[round-mode=places,round-precision=1]{\fidGradNormmean}~FID with zero
recall (Table~\ref{tab:full}), a failure of the balancer at one setting
rather than evidence about balancing. The three-level arm is unstable across
seeds (\num[round-mode=places,round-precision=1]{\fidMultilevelmean}~$\pm$~%
\num[round-mode=places,round-precision=1]{\fidMultilevelsd}~FID); its
per-band averaging weights a level-$j$ detail coefficient $4^{j-1}$ times a
level-1 one (Section~\ref{sec:family-parseval}), so the instability cannot be
attributed to depth alone.

% Generated by scripts/make_tables.py --- do not edit by hand.
% Regenerate after every sweep; the paper reads these directly.

\begin{table}[tbp]
\centering
\footnotesize
\setlength{\tabcolsep}{2.5pt}
\caption{The method extensions, CIFAR-10, epoch 199, DDIM-100; $n$ is the number of seeds. No extension normalises its band weights. The per-sub-band, spatially adaptive and lifting arms keep AWWL's rational weighting with every detail band at the full $w_{\mathrm{det}}$, the summed form; the multi-level arm divides $w_{\mathrm{det}}$ among its nine detail bands; the learned-weighting, GradNorm and all-axes arms set the band weights themselves. Every extension is compared with the unnormalised AWWL with detail bands summed (second row, the parent of the summed-form arms), not with the normalised variants of Table~\ref{tab:confound}. $\Delta$ is the paired-by-seed FID difference with its 95\% $t$-interval, over the seeds both arms share; $p$ is the two-sided paired $t$-test, raw and Holm-corrected: against MSE over the 21-arm family of Table~\ref{tab:compare-fid}, against the parent over the seven extensions. Bold: $p_{\mathrm{Holm}} < 0.05$ within the stated family. Axes turned on: Per-sub-band schedules: separate schedule per detail band; Spatially adaptive: weights vary over position; Multi-level (3 levels): three-level decomposition; Learned weighting: $\sigma$-conditioned learned weights; GradNorm balancing: GradNorm sets the weights; Learned basis (lifting): learned lifting basis; All axes combined: learned weights, spatial, lifting.}
\label{tab:extensions}
\begin{tabular}{lrrrrrrr}
\toprule
 & & \multicolumn{3}{c}{Against MSE} & \multicolumn{3}{c}{Against the parent} \\
\cmidrule(lr){3-5}\cmidrule(lr){6-8}
Arm & $n$ & $\Delta$ [95\% CI] & $p$ & $p_{\mathrm{Holm}}$ & $\Delta$ [95\% CI] & $p$ & $p_{\mathrm{Holm}}$ \\
\midrule
\multicolumn{8}{l}{\emph{Reference}} \\
AWWL & 5 & $+0.44$\;$[-0.15, 1.03]$ & $0.11$ & $0.32$ &  &  &  \\
AWWL, detail bands summed & 5 & $+2.63$\;$[2.22, 3.04]$ & $<10^{-4}$ & $\mathbf{0.0008}$ & \emph{parent} &  &  \\
\midrule
Per-sub-band schedules & 5 & $+2.61$\;$[2.11, 3.11]$ & $0.0001$ & $\mathbf{0.0017}$ & $-0.02$\;$[-0.42, 0.39]$ & $0.92$ & $0.92$ \\
Spatially adaptive & 5 & $+4.65$\;$[4.10, 5.19]$ & $<10^{-4}$ & $\mathbf{0.0003}$ & $+2.02$\;$[1.70, 2.34]$ & $<10^{-4}$ & $\mathbf{0.0003}$ \\
Multi-level (3 levels) & 5 & $+88.9$\;$[50.4, 127.3]$ & $0.0030$ & $\mathbf{0.021}$ & $+86.2$\;$[47.9, 124.6]$ & $0.0034$ & $\mathbf{0.010}$ \\
Learned weighting & 5 & $+6.66$\;$[6.12, 7.20]$ & $<10^{-4}$ & $\mathbf{<10^{-4}}$ & $+4.03$\;$[3.64, 4.42]$ & $<10^{-4}$ & $\mathbf{<10^{-4}}$ \\
GradNorm balancing & 5 & $+330$\;$[296, 365]$ & $<10^{-4}$ & $\mathbf{0.0002}$ & $+328$\;$[293, 363]$ & $<10^{-4}$ & $\mathbf{<10^{-4}}$ \\
Learned basis (lifting) & 3 & $+6.13$\;$[4.99, 7.28]$ & $0.0019$ & $\mathbf{0.015}$ & $+3.36$\;$[2.37, 4.34]$ & $0.0046$ & $\mathbf{0.010}$ \\
All axes combined & 3 & $+12.5$\;$[11.5, 13.5]$ & $0.0004$ & $\mathbf{0.0039}$ & $+9.76$\;$[9.06, 10.45]$ & $0.0003$ & $\mathbf{0.0011}$ \\
\bottomrule
\end{tabular}
\end{table}


\FloatBarrier
}{%
  \todo{sections/supp\_audit.tex has not been written yet.}}
\FloatBarrier
\IfFileExists{sections/supp_mechanism.tex}{% Supplementary material for Section~\ref{sec:mechanism} (Does AWWL change the
% spectrum?). Input under Appendix D (Supplementary results); starts at
% subsection level. Moved floats keep their labels: tab:bands-ddpm,
% fig:compare.

\subsection{The spectrum under DDPM-1000, and paired samples}
\label{app:supp-mechanism}

Table~\ref{tab:bands-ddpm} repeats Table~\ref{tab:bands} for ancestral
DDPM-1000 sampling of the same weights (Section~\ref{sec:mechanism}). Under
this sampler the baseline's deficit is smaller in every band, the static
wavelet arm stays close to the MSE level, as under DDIM-100 ($\corrHighStaticddpm$~dB
$[\corrHighStaticddpmLo, \corrHighStaticddpmHi]$,
$p = \corrHighStaticddpmP$), and AWWL's high-band change is a small shift
towards a larger deficit ($\corrHighAWWLddpm$~dB
$[\corrHighAWWLddpmLo, \corrHighAWWLddpmHi]$, $p = \corrHighAWWLddpmP$).

\IfFileExists{tables/bands_ddpm1000.tex}{% Generated by scripts/make_band_table.py --- do not edit by hand.
% Recomputed from the sampled images (all samples, all reference images).

\begin{table}[tbp]
\centering
\footnotesize
\setlength{\tabcolsep}{3pt}
\caption{Signed deviation of the samples' radial luminance spectrum from the reference set's, in dB, at epoch 199 (DDPM-1000); negative means too little energy. Every sample of each run and all \num{50000} CIFAR-10 training images enter the profiles. Bands are thirds of the integer radial bins: radii 0--6 (low), 7--14 (mid) and 15--22 (high), in cycles per image, where the axial Nyquist radius is 16. A band's value is the unweighted mean of its bins' deviations, so each radius counts once whatever its number of coefficients. Mean $\pm$ sample standard deviation over five seeds. The last two columns pair each seed with its own MSE run in the high band: the difference of deviations (positive means more high-frequency energy than MSE), its 95\% $t$-interval and the two-sided paired $t$-test $p$, uncorrected.}
\label{tab:bands-ddpm}
\begin{tabular}{lrrrrr}
\toprule
Configuration & Low & Mid & High & $\Delta$ high [95\% CI] & $p$ \\
\midrule
MSE & $-0.14 \pm 0.05$ & $-0.42 \pm 0.04$ & $-1.30 \pm 0.04$ & --- & --- \\
AWWL & $-0.13 \pm 0.04$ & $-0.41 \pm 0.04$ & $-1.34 \pm 0.03$ & $-0.04$\;$[-0.06, -0.03]$ & $0.0021$ \\
Static wavelet ($p = 0$) & $-0.14 \pm 0.04$ & $-0.42 \pm 0.04$ & $-1.32 \pm 0.05$ & $-0.02$\;$[-0.05, 0.004]$ & $0.074$ \\
\bottomrule
\end{tabular}
\end{table}
}{%
  \begin{table}[tbp]\centering
  \fbox{\textbf{tables/bands\_ddpm1000.tex not built}}
  \caption{Not built.}\label{tab:bands-ddpm}
  \end{table}}

Figure~\ref{fig:compare} pairs uncurated samples of MSE, the static wavelet
arm and AWWL drawn from the same initial noise. The difference row shows small
differences spread over the whole image and larger ones concentrated at object
boundaries and in small structures. Across the three sample rows, no
systematic difference is apparent on inspection at print size.

\guardedfigure{1.0}{compare.png}{Uncurated samples (indices 0--7) of training
seed 1 at epoch 199, drawn with DDIM-100 ($\eta = 0$). All rows of a column
start from the same initial noise, and DDIM with $\eta = 0$ is deterministic
given that noise, so the images in a column differ only through the training
objective. The last row is $\lvert \mathrm{AWWL} - \mathrm{MSE} \rvert$,
averaged over the RGB channels and amplified eight-fold (darker is larger;
differences of 32 or more grey levels print black).}{fig:compare}
}{%
  \todo{sections/supp\_mechanism.tex has not been written yet.}}
\FloatBarrier
\IfFileExists{sections/supp_confound.tex}{% Supplementary results for Section~\ref{sec:confound}: the normalisation
% control built on the summed-detail variant. Input under the appendix section
% that carries \label{app:supp}, so this file starts at subsection level.

\subsection{A normalisation control on the summed-detail shape}
\label{app:supp-confound}

A further control is built on the summed-detail variant of
Eq.~\eqref{eq:summed}, whose per-band weights are
$(w_{\mathrm{LL}}, w_{\det}, w_{\det}, w_{\det})$, and fixes
$\sum_b \omega_b^2$ at \levelSqFixedSum, that variant's mean over the training
timesteps (Table~\ref{tab:confound-summed}). Because it differs from the
controls of Table~\ref{tab:confound} in band shape as well as in
normalisation, it is compared only with its parent, the unnormalised
summed-detail arm. Against that parent, the control changes FID by
$\dSqFixedSumVsEqSeven$ $[\dSqFixedSumVsEqSevenLo, \dSqFixedSumVsEqSevenHi]$
(raw $p = \dSqFixedSumVsEqSevenP$, $p_{\mathrm{Holm}} = \dSqFixedSumVsEqSevenHolm$
over the \nFamilyConfound{} confound contrasts), although it remains
worse than MSE ($+\dSqFixedSum$~FID $[\dSqFixedSumLo, \dSqFixedSumHi]$,
$p_{\mathrm{Holm}} = \dSqFixedSumHolm$). For this shape, whose total falls
from $3(1 - \alpha)$ to $\alpha$, holding $\sum_b \omega_b^2$ constant
therefore improves FID relative to the parent, whereas for the released,
detail-averaged shape the same normalisation raises FID relative to AWWL
(Section~\ref{sec:confound-results}). The effect of normalising the timestep
weight depends on the band shape it accompanies, and the controls support no
general statement that timestep reweighting helps or harms. As for the main
controls, the effective timestep weight $\Lambda(t)$ is not held fixed
(Proposition~\ref{prop:confound}(b)).

% Generated by scripts/make_tables.py --- do not edit by hand.
% Regenerate after every sweep; the paper reads these directly.

\begin{table}[tbp]
\centering
\footnotesize
\setlength{\tabcolsep}{3pt}
\caption{Normalisation control on the summed-detail variant of AWWL (each detail band carries $w_{\det}$), CIFAR-10, epoch 199, DDIM-100. The band weights are $\omega_b = (w_{\mathrm{LL}}, w_{\det}, w_{\det}, w_{\det})$; the control holds $\sum_b \omega_b^2$ at the variant's own mean over the training timesteps. $\Delta$ is the paired-by-seed FID difference with its 95\% $t$-interval over the $n$ seeds; $p$ is the two-sided paired $t$-test, raw and Holm-corrected. Against MSE the Holm family is the 21-arm family of Table~\ref{tab:compare-fid}; against the reference it is the four confound contrasts of Tables~\ref{tab:confound} and~\ref{tab:confound-summed}. Levels read from the run configurations: AWWL, detail bands summed, $\sum_b \omega_b^2$ fixed: $\sum_b \omega_b^2 = 0.2454$. Bold: $p_{\mathrm{Holm}} < 0.05$ within the stated family.}
\label{tab:confound-summed}
\begin{tabular}{>{\raggedright\arraybackslash\hangindent=1em\hangafter=1}p{3.3cm}rrrrrrr}
\toprule
 & & \multicolumn{3}{c}{Against MSE} & \multicolumn{3}{c}{Against its parent (first row)} \\
\cmidrule(lr){3-5}\cmidrule(lr){6-8}
Arm & $n$ & $\Delta$ [95\% CI] & $p$ & $p_{\mathrm{Holm}}$ & $\Delta$ [95\% CI] & $p$ & $p_{\mathrm{Holm}}$ \\
\midrule
AWWL, detail bands summed & 5 & $+2.63$\;$[2.22, 3.04]$ & $<10^{-4}$ & $\mathbf{0.0008}$ & \emph{reference} &  &  \\
AWWL, detail bands summed, $\sum_b \omega_b^2$ fixed & 5 & $+1.75$\;$[1.19, 2.32]$ & $0.0010$ & $\mathbf{0.0090}$ & $-0.87$\;$[-1.16, -0.58]$ & $0.0011$ & $\mathbf{0.0045}$ \\
\bottomrule
\end{tabular}
\end{table}

}{%
  \todo{sections/supp\_confound.tex has not been written yet.}}
% No barrier here: the summed-shape table may float to the top of the next
% page, beside the DreamBooth text, instead of leaving a part-empty page.
\IfFileExists{sections/supp_generality.tex}{% Supplementary results for Section~\ref{sec:generality}. Input under
% Appendix D (Supplementary results, \label{app:supp}); starts at subsection
% level, as the other appendix inputs do. It carries tab:dreambooth,
% sec:restoration, fig:restoration and tab:restoration, so the Gaussian-denoising
% subsection of appendix.tex (\label{app:tables-restoration}, which also inputs
% tables/restoration) must be removed when this file is input.

\subsection{DreamBooth: means and ranks}
\label{app:supp-dreambooth}

Table~\ref{tab:dreambooth} gives the mean, the rank and the Holm-corrected
$p$-value of each objective on both metrics of the DreamBooth experiment of
Section~\ref{sec:dreambooth}; Fig.~\ref{fig:db-tradeoff} plots the same runs
as paired differences from MSE. The two rankings are nearly the reverse of
each other, so the objectives trade prompt agreement against subject
fidelity, and none is better than MSE on both metrics. For the one difference
that survives correction, raw SNR weighting on CLIP score, the Wilcoxon
signed-rank test gives $p = \dbWilcoxonClipRawSNR$, the smallest value that
test can return with five pairs.

\IfFileExists{tables/dreambooth.tex}{% Generated by scripts/make_tables.py --- do not edit by hand.
% Regenerate after every sweep; the paper reads these directly.

\begin{table}[tbp]
\centering
\small
\caption{DreamBooth fine-tuning of Stable Diffusion~1.5 on one subject, five seeds per objective, \num{60} images per run. Rank is over the ten objectives, the baseline included, on each metric; $p$ is the two-sided paired-by-seed $t$-test against MSE, Holm-corrected over the nine objectives within each metric. Bold: $p_{\mathrm{Holm}} < 0.05$ within the stated family. The published claim is that AWWL is the only method in the top two on both metrics.}
\label{tab:dreambooth}
\begin{tabular}{lrrrrrr}
\toprule
 & \multicolumn{3}{c}{CLIP score $\uparrow$} & \multicolumn{3}{c}{Subject similarity $\uparrow$} \\
\cmidrule(lr){2-4}\cmidrule(lr){5-7}
Objective & Mean & Rank & $p_{\mathrm{Holm}}$ & Mean & Rank & $p_{\mathrm{Holm}}$ \\
\midrule
Raw SNR weighting & $0.3522$ & 1 & $\mathbf{0.0057}$ & $0.7989$ & 10 & $0.085$ \\
Learned weighting & $0.3199$ & 2 & $0.47$ & $0.8388$ & 9 & $1.00$ \\
MSE & $0.3111$ & 3 & --- & $0.8429$ & 7 & --- \\
Huber & $0.3103$ & 4 & $0.90$ & $0.8389$ & 8 & $1.00$ \\
Min-SNR-$\gamma$ & $0.3092$ & 5 & $0.90$ & $0.8493$ & 6 & $1.00$ \\
Static wavelet ($p = 0$) & $0.3049$ & 6 & $0.47$ & $0.8540$ & 5 & $1.00$ \\
AWWL, $\sum_b \omega_b = 1$ & $0.3008$ & 7 & $0.098$ & $0.8635$ & 2 & $0.40$ \\
AWWL & $0.3003$ & 8 & $0.24$ & $0.8721$ & 1 & $0.40$ \\
Charbonnier & $0.2985$ & 9 & $0.10$ & $0.8564$ & 4 & $1.00$ \\
$L_1$ & $0.2967$ & 10 & $0.10$ & $0.8634$ & 3 & $1.00$ \\
\bottomrule
\end{tabular}
\end{table}
}{%
  \begin{table}[tbp]\centering
  \fbox{\parbox{0.8\linewidth}{\centering\vspace{4mm}
  \textbf{tables/dreambooth.tex not built}\\[1mm]
  run \texttt{python scripts/make\_tables.py}\vspace{4mm}}}
  \caption{DreamBooth, five seeds per objective.}\label{tab:dreambooth}
  \end{table}}

\subsection{Gaussian denoising}
\label{sec:restoration}

A frequency-aware objective might pay where the target has image structure and
the metrics score the output directly. We train the objectives as direct
denoisers (not diffusion models) on CIFAR-10 with $\sigma \in [0.05, 0.5]$
drawn per image, three seeds of 100 epochs, scored on 500 test images by
PSNR, SSIM \citep{wang2004image}, LPIPS \citep{zhang2018unreasonable} and a
NIQE-feature distance \citep{mittal2013making} whose pristine model is fitted
on the clean test images (Appendix~\ref{app:setup-restoration}). The three wavelet
arms (AWWL; AWWL, $\sum_b \omega_b = 1$; static wavelet, $p = 0$) apply AWWL's
weighting with the degradation level as $\sigma$ at $\alpha = 1/2$, with
$p = 1$ except in the static wavelet arm. The total weight $S$
of Eq.~\eqref{eq:total} is then constant, so the weighting has no
residual-independent noise-level component (Proposition~\ref{prop:confound}),
and normalising the weights to sum to 1 only doubles them.

\guardedfigure{1.0}{restoration_null.png}{%
  Gaussian denoising on CIFAR-10 ($\sigma \in [0.05, 0.5]$, 500 test images,
  degradation shared across arms): paired difference from MSE in (a) PSNR,
  (b) SSIM, (c) LPIPS and (d) the NIQE-feature distance ($\Delta$NIQE), mean
  of three seeds with paired 95\% $t$-intervals. Hollow: no difference
  survives Holm correction within any metric; two intervals exclude zero
  before correction (Charbonnier on PSNR; AWWL, $\sum_b \omega_b = 1$, on the
  NIQE-feature distance). Wavelet arms use $\alpha = 1/2$ and $p = 1$, except
  the static wavelet arm ($p = 0$).}{fig:restoration}

\IfFileExists{tables/restoration.tex}{% Generated by scripts/make_tables.py --- do not edit by hand.
% Regenerate after every sweep; the paper reads these directly.

\begin{table}[tbp]
\centering
\small
\caption{Gaussian denoising on CIFAR-10, three seeds, epoch \num{99}; means over seeds. The AWWL arms use $\alpha = 0.5$. Each arm is paired by seed with MSE (two-sided paired $t$-test). Paired differences with 95\% intervals are plotted in Fig.~\ref{fig:restoration}; the smallest raw $p$ is $0.027$ (AWWL, $\sum_b \omega_b = 1$, NIQE-feature distance), $0.13$ after Holm correction over the five arms within that metric. NIQE-feature distance is the distance of the NIQE features from a pristine model fitted on the clean test images. The degradation is drawn once and reused across every arm.}
\label{tab:restoration}
\begin{tabular}{lrrrr}
\toprule
Objective & PSNR $\uparrow$ & SSIM $\uparrow$ & LPIPS $\downarrow$ & NIQE-feat.\ dist.\ $\downarrow$ \\
\midrule
MSE & $28.802$ & $0.8698$ & $0.01613$ & $1.5195$ \\
AWWL & $28.776$ & $0.8693$ & $0.01592$ & $1.4942$ \\
Static wavelet ($p = 0$) & $28.773$ & $0.8692$ & $0.01579$ & $1.5065$ \\
AWWL, $\sum_b \omega_b = 1$ & $28.779$ & $0.8698$ & $0.01597$ & $1.4960$ \\
$L_1$ & $28.801$ & $0.8703$ & $0.01680$ & $1.5618$ \\
Charbonnier & $28.798$ & $0.8699$ & $0.01636$ & $1.5389$ \\
\bottomrule
\end{tabular}
\end{table}
}{%
  \begin{table}[tbp]\centering
  \fbox{\textbf{tables/restoration.tex not built}}
  \caption{Not built.}\label{tab:restoration}
  \end{table}}

No arm differs from MSE after Holm correction (Fig.~\ref{fig:restoration},
Table~\ref{tab:restoration}); the smallest raw $p$ over the four metrics is
\restMinP{} (\restMinPHolm{} after correction). The wavelet arms lean towards
lower PSNR and better LPIPS and NIQE-feature distance, the direction of the
perception--distortion trade-off \citep{blau2018perception}, by amounts that
three seeds do not resolve. The 95\% interval of AWWL on PSNR,
$[\restDPSNRAWWLLo, \restDPSNRAWWLHi]$~dB, excludes a gain above
\restDPSNRAWWLHi~dB. On the half of the test images with the lowest noise
levels, all three wavelet arms have lower PSNR than MSE, with intervals that
exclude zero before correction; for AWWL the difference is
$\restDPSNRlowAWWL$~dB $[\restDPSNRlowAWWLLo, \restDPSNRlowAWWLHi]$ (raw
$p = \restDPSNRlowAWWLP$, $p_{\mathrm{Holm}} = \restDPSNRlowAWWLHolm$). This
task was not priced: whether LPIPS and the NIQE-feature distance register the
high band at 32 pixels (precondition P4) is untested.
}{%
  \todo{sections/supp\_generality.tex has not been written yet.}}
\FloatBarrier
